\documentclass[11pt,a4paper]{article}
\usepackage[T1]{fontenc}
\usepackage[utf8]{inputenc}
\usepackage{libertinus}
\usepackage{microtype}
\usepackage[a4paper,margin=26mm,headheight=15pt]{geometry}
\usepackage{amsmath,amssymb,amsthm,mathtools,bm,mathrsfs}
\usepackage{booktabs,longtable,tabularx,array}
\usepackage{xcolor}
\usepackage{enumitem}
\usepackage{fancyhdr}
\usepackage{xurl}
\usepackage{hyperref}
\usepackage[nameinlink,noabbrev,capitalise]{cleveref}
\definecolor{deepolive}{RGB}{67,79,45}
\definecolor{darkteal}{RGB}{31,83,84}
\hypersetup{colorlinks=true,linkcolor=deepolive,citecolor=darkteal,urlcolor=darkteal,
 pdftitle={Gaussian Coefficients in Combinatorial Coefficient Calculus},
 pdfsubject={q-binomial identities, Newton bases, Stirling and Eulerian transforms, Bell polynomials, inversion, and asymptotics}}
\pagestyle{fancy}\fancyhf{}
\fancyhead[L]{\small Gaussian coefficient calculus}
\fancyhead[R]{\small\thepage}
\renewcommand{\headrulewidth}{0.4pt}
\fancypagestyle{plain}{\fancyhf{}\fancyfoot[C]{\thepage}\renewcommand{\headrulewidth}{0pt}}
\setlength{\emergencystretch}{2em}
\setlist{itemsep=2pt,topsep=4pt}
\allowdisplaybreaks[2]\raggedbottom
\numberwithin{equation}{section}
\newtheorem{theorem}{Theorem}[section]
\newtheorem{lemma}[theorem]{Lemma}
\newtheorem{proposition}[theorem]{Proposition}
\newtheorem{corollary}[theorem]{Corollary}
\theoremstyle{definition}
\newtheorem{definition}[theorem]{Definition}
\newtheorem{example}[theorem]{Example}
\theoremstyle{remark}
\newtheorem{remark}[theorem]{Remark}
\newtheorem{warning}[theorem]{Warning}
\newcommand{\N}{\mathbb N}
\newcommand{\Z}{\mathbb Z}
\newcommand{\Q}{\mathbb Q}
\newcommand{\C}{\mathbb C}
\newcommand{\R}{\mathbb R}
\newcommand{\G}[2]{\genfrac{[}{]}{0pt}{}{#1}{#2}_{q}}
\newcommand{\Gb}[3]{\genfrac{[}{]}{0pt}{}{#1}{#2}_{#3}}
\newcommand{\qnum}[1]{[#1]_q}
\newcommand{\qfac}[1]{[#1]_q!}
\newcommand{\qp}[2]{(#1;q)_{#2}}
\newcommand{\coef}[2]{[#1^{#2}]}
\newcommand{\ct}{\operatorname{CT}}
\newcommand{\Res}{\operatorname{Res}}
\newcommand{\Dq}{D_q}
\newcommand{\Li}{\operatorname{Li}}
\newcommand{\des}{\operatorname{des}}
\newcommand{\maj}{\operatorname{maj}}
\newcommand{\inv}{\operatorname{inv}}
\newcommand{\Sg}[2]{\left\{\begin{matrix}#1\\#2\end{matrix}\right\}}
\newcommand{\fall}[2]{#1^{\underline{#2}}}
\newcommand{\dd}{\,\mathrm d}
\newcommand{\ee}{\mathrm e}
\newcommand{\id}{\mathrm{id}}
\newcommand{\one}{\mathbf 1}
\newcommand{\Qn}{Q_n}
\newcommand{\Bcal}{\mathcal B}
\newcommand{\Y}{\mathcal Y}
\newcommand{\Poch}{\mathcal P}
\newcommand{\srcpath}{\nolinkurl{Analysis/FabiusFunction/docs/semi-formalized-research-frontiers/drafts/combinatorial-coefficient-calculus/Combinatorial_Coefficient_Calculus}}
\title{\Huge\bfseries Gaussian Coefficients in\\[2mm]
Combinatorial Coefficient Calculus\\[5mm]
\Large A $q$-binomial extension through Newton bases,\\
Stirling--Eulerian transforms, Bell polynomials,\\
series inversion, and asymptotic expansions}
\author{A self-contained research companion to\\\emph{Combinatorial Coefficient Calculus}}
\date{4 September 2026}
\begin{document}
\maketitle
\begin{abstract}
Gaussian, or $q$-binomial, coefficients are more than weighted versions of
ordinary binomial coefficients: they are the connection coefficients between
monomials and factorial polynomials on a geometric grid. This observation
places them directly inside the basis-change and coefficient-extraction
framework of \emph{Combinatorial Coefficient Calculus}. We develop that
extension systematically. Finite products give convolution and inversion;
affine changes of geometric nodes give explicit $q$-Stirling and colored-node
formulae; weighted sorting gives $q$-Worpitzky and Stirling--Eulerian
transforms. Ordinary Bell polynomials organize powers of $q$-products,
parameter derivatives, local expansions at $q=1$ and at roots of unity, and
closed coefficients for compositional inverses. Jackson differentiation is
handled separately from ordinary composition, including a divided-difference
chain rule that retains the necessary geometric shifts. Three asymptotic
regimes are distinguished, and an application gives exact Gaussian
extrapolation of moments of geometric uniform convolutions, including the
Rvachev half-base case. All principal identities are proved; an accompanying
exact-arithmetic program checks finite instances independently. The
contribution is an integrated extension and a collection of explicit
coefficient mechanisms, not a claim that the classical identities are new.
\end{abstract}
\thispagestyle{plain}
\clearpage
\setcounter{tocdepth}{1}
\hypersetup{bookmarksdepth=2}
\tableofcontents
\clearpage

\section{Scope, conventions, and the connection to the parent article}
\label{sec:scope}
The parent monograph \cite{parent} organizes combinatorics around polynomial
basis changes, ordinary and exponential generating functions, Bell-polynomial
composition, and Lagrange--B\"urmann inversion. The present article extends
those mechanisms rather than replacing the binomial sign in every formula
by a Gaussian bracket. This distinction is essential: there are several
inequivalent objects called $q$-Stirling, $q$-Bell, or $q$-Catalan numbers.
The defining identities below, not the names, fix our conventions.

The parent source was consulted on the live \texttt{main} branch on
4 September 2026. The sections on formal series, Newton expansion, Bell
polynomials, composition, Lagrange inversion, and Riordan arrays determine
our points of contact. Since its chapter numbering can change, we refer to
these topics rather than frozen page numbers. The statements here are not presented as Lean-checked results. The repository's broader
Fabius--Rvachev exposition already contains substantial $q$-calculus
\cite{fabiusrepo}; the classical ingredients repeated here serve as a
self-contained foundation for the connections developed below.

\subsection{Algebraic and analytic settings}
For algebraic identities, $q$ is initially an indeterminate. Polynomial
identities live in $\Z[q]$; divisions are performed in $\Q(q)$ or in a
specified extension. Formal power series in $z$ are taken over a
commutative $\Q(q)$-algebra. A product $\prod_{j\ge0}(1-aq^jz)$ is not,
without an additional topology or a coefficient definition, a product in
$\Q(q)[[z]]$. We first define its coefficients algebraically or work
analytically with $|q|<1$, and only then use infinite-product notation.

The Gaussian polynomial is evaluated at $q=1$, $q=0$, or a root of unity
\emph{after} polynomial cancellation. Expressions containing
$1/(q;q)_n$ or $1/\qfac n$ do not generally survive those specializations.
Unless otherwise stated, $n,k,N$ are nonnegative integers, empty products
are one, impossible finite sums are zero, and $0^0=1$ in coefficient
formulae. An ordinary binomial coefficient with nonnegative upper index is
zero outside its usual range. The default zero convention for Gaussian
coefficients likewise concerns a \emph{nonnegative} upper index; a separate
negative-index convention appears in \cref{sec:infinite}.

\subsection{A map of the extension}
\begin{center}
\small
\begin{tabularx}{\textwidth}{@{}p{0.28\textwidth}X@{}}
\toprule
Parent mechanism & Gaussian extension in this article\\
\midrule
Factorial bases and Stirling inversion & Geometric Newton bases; affine-geometric
nodes; explicit $q$-Stirling transforms (\cref{sec:newton,sec:stirling}).\\
Eulerian and Worpitzky identities & Descent--major-index enumeration and the
Gaussian Stirling--Eulerian bridge (\cref{sec:eulerian}).\\
Bell polynomials and coefficient powers & Geometric power sums, signed alphabets,
and $q$-normalized composition (\cref{sec:bell,sec:composition}).\\
Fa\`a di Bruno calculus & Jackson product rules and a divided-difference chain
rule on geometric nodes (\cref{sec:jackson,sec:chain}).\\
Lagrange and Riordan inversion & Closed inverse coefficients for finite
$q$-products, and Gaussian Pascal arrays (\cref{sec:composition,sec:lagrange}).\\
Bernoulli and asymptotic calculus & All-order local $q$-expansions, root-of-unity
regularization, and a uniform double-scaling expansion
(\cref{sec:derivatives,sec:nearone,sec:roots,sec:asymptotic}).\\
\bottomrule
\end{tabularx}
\end{center}
Three distinctions will recur. A change from $n!$ to $\qfac n$ is a
\emph{normalization of coefficients}; replacing arithmetic nodes by
geometric nodes is a \emph{change of polynomial basis}; replacing the
ordinary derivative by $D_q$ is a \emph{change of operator}. They interact,
but they are not the same construction.

\section{Gaussian polynomials and their finite models}
\label{sec:gaussian}
\subsection{Definitions, Pascal identities, and symmetry}
Set
\begin{equation}
 \qnum n=1+q+\cdots+q^{n-1},\qquad
 \qfac n=\prod_{j=1}^n\qnum j,\qquad
 \qp aN=\prod_{j=0}^{N-1}(1-aq^j).
 \label{eq:qdefs}
\end{equation}
Thus $\qnum0=0$, $\qfac0=1$, and
$\qp qn=(1-q)^n\qfac n$. The notation in this section is standard;
see \cite{dlmf17} for the classical $q$-calculus conventions.

\begin{definition}
For $0\le k\le n$, define
\begin{equation}
 G_{n,k}(q)=\G nk
 =\frac{\qp qn}{\qp qk\qp q{n-k}}
 =\prod_{i=1}^k\frac{1-q^{n-k+i}}{1-q^i}.
 \label{eq:gaussdef}
\end{equation}
For $k<0$ or $k>n\ge0$, set $G_{n,k}(q)=0$.
\end{definition}

\begin{theorem}[Pascal identities and polynomiality]
\label{thm:pascal}
The quotients in \cref{eq:gaussdef} are polynomials in $\Z[q]$. For $n\ge1$ they satisfy
\begin{align}
 \G nk&=\G{n-1}{k-1}+q^k\G{n-1}k,
 \label{eq:pascal1}\\
 \G nk&=\G{n-1}k+q^{n-k}\G{n-1}{k-1}.
 \label{eq:pascal2}
\end{align}
For $0\le k\le n$,
\begin{equation}
 \G nk=\G n{n-k},\qquad
 G_{n,k}(q^{-1})=q^{-k(n-k)}G_{n,k}(q),\qquad
 G_{n,k}(1)=\binom nk.
 \label{eq:symmetries}
\end{equation}
Their coefficients are nonnegative integers, and their degree is $k(n-k)$.
\end{theorem}
\begin{proof}
For $0<k<n$, put the two terms in \cref{eq:pascal1} over the
common denominator $\qp qk\qp q{n-k}$. Their numerator is
$\qp q{n-1}\{1-q^k+q^k(1-q^{n-k})\}=\qp qn$.
The boundary cases are immediate. This proves polynomiality and
nonnegativity inductively, starting with $G_{0,0}=1$. Symmetry follows from
the factorial quotient, and gives \cref{eq:pascal2}. Replacing each factor
$1-q^{-j}$ by $-q^{-j}(1-q^j)$ gives the reciprocity law.
The factors $(1-q^a)/(1-q)$ tend to $a$, giving the ordinary binomial
limit. The degree assertion follows from the product, or from the
weighted model in the next proof, whose minimum and maximum weights are
both attained.
\end{proof}

\subsection{Subsets, rectangle partitions, and words}
\begin{theorem}[Three polynomial models]
\label{thm:subset}
For $0\le k\le n$,
\begin{align}
 \G nk
 &=\sum_{0\le i_1<\cdots<i_k<n}
       q^{i_1+\cdots+i_k-\binom k2},
 \label{eq:subset}\\
 &=\sum_{0\le b_1\le\cdots\le b_k\le n-k}q^{b_1+\cdots+b_k}.
 \label{eq:rectangle}
\end{align}
Equivalently, it enumerates partitions fitting in a $k\times(n-k)$
rectangle by their size. It also enumerates binary words with $k$ ones
and $n-k$ zeros by the number of pairs consisting of a zero to the left
of a one.
\end{theorem}
\begin{proof}
Let the first sum be $H_{n,k}$. Subsets not containing $n-1$ contribute
$H_{n-1,k}$; deleting $n-1$ from the others decreases the exponent by
$(n-1)-(k-1)=n-k$. Thus $H$ satisfies \cref{eq:pascal2} with the same
boundary values as $G$. The substitution $b_j=i_j-(j-1)$ gives the second
sum. Reversing the $b_j$ gives a partition inside the rectangle.
In the associated binary word, the $j$th one has exactly
$i_j-(j-1)$ zeros before it. Summing gives the stated word statistic.
\end{proof}
Reversal of the word changes that statistic to $k(n-k)$ minus itself,
which explains reciprocity without division. For example,
\begin{equation}
 \G42=1+q+2q^2+q^3+q^4.
 \label{eq:G42}
\end{equation}
The two objects of weight $2$ are the rectangle partitions $(2,0)$ and
$(1,1)$.

\subsection{Multinomials, finite fields, and flags}
For $n_1+\cdots+n_r=n$, write
\begin{equation}
 \Gb n{n_1,\ldots,n_r}q
 =\frac{\qfac n}{\prod_{i=1}^r\qfac{n_i}}
 =\prod_{i=1}^r\G{n-\sum_{j<i}n_j}{n_i}.
 \label{eq:qmultinomial}
\end{equation}
The right side proves polynomiality. Assigning the positions of letter
$1$, then letter $2$, and so on shows that it enumerates words with these
multiplicities by the usual inversion statistic; reversal also gives the
complementary convention.

\begin{proposition}[Subspace and flag counts]
\label{prop:flags}
If $Q$ is a prime power, $G_{n,k}(Q)$ counts the $k$-dimensional
subspaces of $\mathbb F_Q^n$. The multinomial in
\cref{eq:qmultinomial}, evaluated at $q=Q$, counts flags whose successive
dimension increments are $n_1,\ldots,n_r$.
\end{proposition}
\begin{proof}
There are $(Q^n-1)(Q^n-Q)\cdots(Q^n-Q^{k-1})$ ordered independent
$k$-tuples in $\mathbb F_Q^n$. Each $k$-subspace has
$(Q^k-1)(Q^k-Q)\cdots(Q^k-Q^{k-1})$ ordered bases. Dividing and canceling
the powers of $Q$ gives \cref{eq:gaussdef}. Choose each next subspace
in the quotient by the preceding one to obtain the flag product.
\end{proof}
The indeterminate $q$ is not itself a field size. The finite-field
interpretation applies only at prime powers; the polynomial identities
hold independently of it.

\section{Finite products, convolutions, and Gaussian inversion}
\label{sec:products}
\subsection{The two finite coefficient theorems}
\begin{theorem}[Finite $q$-binomial theorem]
\label{thm:finite}
For $N\ge0$,
\begin{equation}
 \qp zN=\sum_{k=0}^N(-1)^kq^{\binom k2}\G Nk z^k.
 \label{eq:finite}
\end{equation}
For $N\ge1$, its reciprocal has the formal expansion
\begin{equation}
 \frac1{\qp zN}=\sum_{m\ge0}\G{N+m-1}m z^m.
 \label{eq:finite-recip}
\end{equation}
\end{theorem}
\begin{proof}
In the product defining $\qp zN$, choose the term $-zq^j$ at a subset of
$k$ positions; \cref{eq:subset} proves \cref{eq:finite}.
In the reciprocal product choose a nonnegative multiplicity of each
$q^jz$. Its degree-$m$ coefficient is
\[
 \sum_{0\le i_1\le\cdots\le i_m<N}q^{i_1+\cdots+i_m}.
\]
The substitution $b_j=i_j+j-1$ gives a strictly increasing $m$-subset of
$\{0,\ldots,N+m-2\}$, with the shift $\binom m2$ removed from its weight.
This is \cref{eq:subset} with upper index $N+m-1$.
\end{proof}

\subsection{Vandermonde identities with all phases visible}
\begin{theorem}[Two-block and multiblock convolution]
\label{thm:vandermonde}
For $m,n,k\ge0$,
\begin{equation}
 \G{m+n}k
 =\sum_{j=0}^k q^{(m-j)(k-j)}\G mj\G n{k-j}.
 \label{eq:vandermonde}
\end{equation}
Terms with invalid indices are omitted, so the displayed powers are
nonnegative on the contributing range. More generally, if
$N=N_1+\cdots+N_r$, then
\begin{equation}
 \G Nk=
 \sum_{\substack{k_1+\cdots+k_r=k\\0\le k_i\le N_i}}
 q^{\sum_{i<j}(N_i-k_i)k_j}\prod_{i=1}^r\G{N_i}{k_i}.
 \label{eq:multivandermonde}
\end{equation}
For $m,n\ge1$, a reciprocal-product version is
\begin{equation}
 \G{m+n+k-1}k
 =\sum_{j=0}^k q^{m(k-j)}
   \G{m+j-1}j\G{n+k-j-1}{k-j}.
 \label{eq:negative-vandermonde}
\end{equation}
\end{theorem}
\begin{proof}
Split the finite product as
$\qp z{m+n}=\qp zm\qp{q^mz}n$ and compare $z^k$ in
\cref{eq:finite}. The exponent left after canceling
$q^{\binom k2}$ is
$\binom j2+\binom{k-j}2+m(k-j)-\binom k2=(m-j)(k-j)$.
For several blocks the same subtraction is
\[
 \sum_i\binom{k_i}2+
 \sum_i k_i\sum_{h<i}N_h-\binom k2
 =\sum_{h<i}(N_h-k_h)k_i.
\]
Finally split the reciprocal product and use
\cref{eq:finite-recip}. This time no triangular powers are canceled;
the second block contributes exactly $q^{m(k-j)}$.
\end{proof}
The phase in \cref{eq:vandermonde} is not optional. It records crossings
between two blocks, or, in the binary model, zeros in the first block
lying before ones in the second.

\subsection{Annihilation and inversion}
\begin{corollary}[Gaussian annihilator]
\label{cor:annihilator}
For an indeterminate $x$,
\begin{equation}
 \sum_{j=0}^N(-1)^jq^{\binom j2}\G Nj x^j=\qp xN.
 \label{eq:annihilator}
\end{equation}
In particular the sum at $x=1$ is zero for $N>0$. If $q\ne0$ and
$x=q^{-r}$ with $0\le r<N$, it is again zero.
\end{corollary}
\begin{proof}
This is \cref{eq:finite}; in the two specializations one factor of the
product is zero.
\end{proof}

\begin{theorem}[Gaussian binomial inversion]
\label{thm:inversion}
For any two sequences with values in a $\Z[q]$-module,
\begin{equation}
 b_n=\sum_{k=0}^n\G nk a_k
 \quad\Longleftrightarrow\quad
 a_n=\sum_{k=0}^n(-1)^{n-k}q^{\binom{n-k}2}\G nk b_k.
 \label{eq:inversion}
\end{equation}
\end{theorem}
\begin{proof}
Substitute the first relation in the right side. The coefficient of
$a_j$ becomes, using factorial cancellation,
\[
 \G nj\sum_{k=j}^n(-1)^{n-k}q^{\binom{n-k}2}\G{n-j}{k-j}.
\]
Put $r=n-k$ and apply \cref{cor:annihilator}. The inner sum is zero
unless $j=n$, when it is one. Both matrices are unit lower triangular,
so this left inverse is also the inverse; alternatively, the same finite
coefficient computation gives the reverse implication.
\end{proof}

\begin{proposition}[M\"obius inversion on subspaces]
If $U\subseteq V$ are finite-dimensional vector spaces over $\mathbb F_Q$
and $d=\dim(V/U)$, the M\"obius function of the subspace lattice is
\begin{equation}
 \mu(U,V)=(-1)^dQ^{\binom d2}.
 \label{eq:mobius}
\end{equation}
Consequently, summation over subspaces of a fixed dimension gives
\cref{eq:inversion} at $q=Q$.
\end{proposition}
\begin{proof}
The defining relation for $\mu$ is
$\sum_{U\subseteq W\subseteq V}\mu(U,W)=\delta_{U,V}$.
There are $G_{d,j}(Q)$ intermediate subspaces with $\dim(W/U)=j$.
The proposed formula satisfies the defining relation by
\cref{cor:annihilator}, and the relation determines $\mu$ uniquely
by increasing interval length.
\end{proof}

\section{Infinite coefficient series and terminating summations}
\label{sec:infinite}
The finite theory is enough for most of this article. The following
infinite series give a convenient algebraic language and a route to
terminating hypergeometric summations; compare \cite{dlmf17,dlmf175,dlmf176}.

\begin{theorem}[The parameter $q$-binomial series]
\label{thm:infinite}
In $\Q(q,a)[[z]]$ define
\begin{equation}
 H_a(z)=\sum_{m\ge0}\frac{\qp am}{\qp qm}z^m.
 \label{eq:Ha}
\end{equation}
It is the unique series with constant term one satisfying
\begin{equation}
 (1-z)H_a(z)=(1-az)H_a(qz).
 \label{eq:Hfunctional}
\end{equation}
When $|q|<1$ and $|z|<1$,
\begin{equation}
 H_a(z)=\frac{\qp{az}\infty}{\qp z\infty}.
 \label{eq:infinite}
\end{equation}
Moreover,
\begin{equation}
 \sum_{j=0}^n\G nj\qp aj\qp b{n-j}a^{n-j}=\qp{ab}n.
 \label{eq:parameter-convolution}
\end{equation}
\end{theorem}
\begin{proof}
Writing $H_a=\sum h_mz^m$ in \cref{eq:Hfunctional} gives
$(1-q^m)h_m=(1-aq^{m-1})h_{m-1}$ for $m\ge1$, which is solved by
the finite product in \cref{eq:Ha}. This proves existence and uniqueness
formally. Analytically the products converge uniformly on compact
$z$-sets away from denominator zeros. Their ratio has constant term one
and satisfies the same equation, proving \cref{eq:infinite}.
The product $H_a(z)H_b(az)$ satisfies \cref{eq:Hfunctional} with $a$
replaced by $ab$: the two intermediate factors $1-az$ cancel. Uniqueness
therefore gives $H_a(z)H_b(az)=H_{ab}(z)$. Compare the coefficient of
$z^n$ and multiply by $\qp qn$.
\end{proof}

In particular, define the \emph{normalized} $q$-exponentials by
\begin{equation}
 e_q(z)=\sum_{n\ge0}\frac{z^n}{\qfac n},\qquad
 E_q(z)=\sum_{n\ge0}\frac{q^{\binom n2}z^n}{\qfac n}.
 \label{eq:qexp}
\end{equation}
The normalization differs from the common convention with $(q;q)_n$
in the denominator by the substitution $z\mapsto(1-q)z$.
For $|q|<1$,
\begin{equation}
 e_q(z)=\frac1{\qp{(1-q)z}\infty},\qquad
 E_q(z)=\qp{-(1-q)z}\infty,\qquad e_q(z)E_q(-z)=1.
 \label{eq:qexp-product}
\end{equation}
The first identity follows from \cref{thm:infinite}. For the second,
expand a finite product and pass coefficientwise to the limit, using
$G_{N,k}(q)\to1/(q;q)_k$. The inverse relation also follows formally
from \cref{cor:annihilator}. The analytic series for $e_q$ is used
inside its disk $|(1-q)z|<1$; $E_q$ is entire in $z$.

\begin{theorem}[The two terminating $q$-Chu--Vandermonde sums]
\label{thm:chu}
For $n\ge0$, as rational identities in generic parameters,
\begin{align}
 \sum_{j=0}^n
 \frac{\qp{q^{-n}}j\qp aj}{\qp qj\qp cj}
       \left(\frac{cq^n}{a}\right)^j
 &=\frac{\qp{c/a}n}{\qp cn},
 \label{eq:chu1}\\
 \sum_{j=0}^n
 \frac{\qp{q^{-n}}j\qp aj}{\qp qj\qp cj}q^j
 &=a^n\frac{\qp{c/a}n}{\qp cn}.
 \label{eq:chu2}
\end{align}
Numerical specialization requires nonzero displayed denominators, or
prior rational cancellation.
\end{theorem}
\begin{proof}
Divide \cref{eq:parameter-convolution} by $a^n\qp bn$ and use
\begin{equation}
 \G nj\frac{\qp b{n-j}}{\qp bn}
 =\frac{\qp{q^{-n}}j}{\qp qj\qp{q^{1-n}/b}j}
       \left(\frac qb\right)^j.
 \label{eq:chu-conversion}
\end{equation}
This equality follows by reversing the $j$ factors canceled from each
product. Set $c=q^{1-n}/b$. Reversing the $n$ factors on the right side
also gives
$\qp{ab}n/(a^n\qp bn)=\qp{c/a}n/\qp cn$, proving
\cref{eq:chu1}. In that identity replace
$(q,a,c)$ by $(q^{-1},a^{-1},c^{-1})$ and use
\[
 (x;q^{-1})_j=(-x)^jq^{-\binom j2}(x^{-1};q)_j.
\]
The transformed summand acquires the net factor $q^j$, and the
transformed right side acquires $a^n$. This proves \cref{eq:chu2}.
\end{proof}

\begin{proposition}[A negative upper index]
For $m\ge1$ and $k\ge0$, define
$\Gb{-m}kq=\prod_{j=0}^{k-1}(1-q^{-m-j})/(q;q)_k$. Then
\begin{equation}
 \Gb{-m}kq=(-1)^kq^{-mk-\binom k2}\G{m+k-1}k.
 \label{eq:negative-top}
\end{equation}
\end{proposition}
\begin{proof}
Factor $-q^{-m-j}$ from the $j$th numerator factor and reverse the
remaining factors. Their quotient by $(q;q)_k$ is the Gaussian
polynomial on the right.
\end{proof}
This limited extension has an unambiguous nonnegative lower index.
Extending \emph{both} indices to all integers requires further
conventions; see the primary treatment \cite{formichella}.

\section{Newton bases: Gaussian coefficients as connection coefficients}
\label{sec:newton}
\subsection{The general node calculus}
Let $a_0,a_1,\ldots$ be elements of a commutative ring and set
\begin{equation}
 P_0^{\boldsymbol a}(x)=1,\qquad
 P_k^{\boldsymbol a}(x)=\prod_{j=0}^{k-1}(x-a_j).
 \label{eq:node-basis}
\end{equation}
Since these polynomials are monic of successive degrees, there are unique
arrays $S_{\boldsymbol a}$ and $s_{\boldsymbol a}$ such that
\begin{equation}
 x^n=\sum_{k=0}^nS_{\boldsymbol a}(n,k)P_k^{\boldsymbol a}(x),\qquad
 P_n^{\boldsymbol a}(x)=\sum_{k=0}^ns_{\boldsymbol a}(n,k)x^k.
 \label{eq:node-connections}
\end{equation}
This common framework for binomial, Stirling, and Gaussian coefficients
is classical; Voigt's short paper \cite{voigt} is particularly close in
spirit to the basis-change program of the parent monograph.

For a finite list $A=(a_0,\ldots,a_k)$, the elementary and complete
homogeneous polynomials are defined by
\[
 e_m(A)=\sum_{i_1<\cdots<i_m}a_{i_1}\cdots a_{i_m},\qquad
 h_m(A)=\sum_{i_1\le\cdots\le i_m}a_{i_1}\cdots a_{i_m}.
\]
Both have value one at $m=0$; $e_m=0$ above the list length.

\begin{theorem}[Universal finite node formulae]
\label{thm:node}
For $0\le k\le n$,
\begin{align}
 S_{\boldsymbol a}(n,k)&=h_{n-k}(a_0,\ldots,a_k),
 \label{eq:node-h}\\
 s_{\boldsymbol a}(n,k)&=(-1)^{n-k}e_{n-k}(a_0,\ldots,a_{n-1}).
 \label{eq:node-e}
\end{align}
The arrays are inverse lower-triangular matrices. In particular,
\begin{equation}
 \sum_{j=k}^ns_{\boldsymbol a}(n,j)S_{\boldsymbol a}(j,k)=\delta_{n,k}.
 \label{eq:node-inverse}
\end{equation}
Their recurrences are
\begin{align}
 S_{\boldsymbol a}(n+1,k)&=S_{\boldsymbol a}(n,k-1)+a_kS_{\boldsymbol a}(n,k),
 \label{eq:node-Srec}\\
 s_{\boldsymbol a}(n+1,k)&=s_{\boldsymbol a}(n,k-1)-a_ns_{\boldsymbol a}(n,k).
 \label{eq:node-srec}
\end{align}
\end{theorem}
\begin{proof}
The identity $xP_k=P_{k+1}+a_kP_k$ gives
\cref{eq:node-Srec}. Separating monomials in $h_m(a_0,\ldots,a_k)$
according as they contain $a_k$ gives the same recurrence for
\cref{eq:node-h}, including the boundary values. Formula
\cref{eq:node-e} follows by expanding the product, and multiplication by
$x-a_n$ gives \cref{eq:node-srec}. Substituting one basis expansion into
the other proves matrix inversion by uniqueness of monomial coefficients.
\end{proof}
All these identities are valid even when nodes coincide. Only the next
representation requires distinct nodes.

\begin{theorem}[Divided differences and a nonrecursive spectral sum]
\label{thm:divdiff}
Over a field, assume $a_0,\ldots,a_k$ are distinct. Then
\begin{equation}
 S_{\boldsymbol a}(n,k)
 =\sum_{j=0}^k\frac{a_j^n}{\prod_{0\le i\le k,\ i\ne j}(a_j-a_i)}.
 \label{eq:node-spectral}
\end{equation}
More generally, the coefficient of $P_k^{\boldsymbol a}$ in the Newton
interpolant to values $f(a_0),\ldots,f(a_M)$ is
\begin{equation}
 f[a_0,\ldots,a_k]
 =\sum_{j=0}^k\frac{f(a_j)}{\prod_{i\ne j}(a_j-a_i)}.
 \label{eq:divided-difference}
\end{equation}
Thus its degree-$M$ interpolating polynomial is
$\sum_{k=0}^Mf[a_0,\ldots,a_k]P_k^{\boldsymbol a}(x)$.
\end{theorem}
\begin{proof}
The Lagrange interpolant on the first $k+1$ nodes is
\[
 L_k(x)=\sum_{j=0}^kf(a_j)
 \frac{\prod_{i\ne j}(x-a_i)}{\prod_{i\ne j}(a_j-a_i)}.
\]
Its leading coefficient is the sum in \cref{eq:divided-difference}.
In a Newton expansion, every term with index greater than $k$ vanishes
on these nodes, and every term with index less than $k$ has smaller
degree. Hence that leading coefficient is precisely the $k$th Newton
coefficient. Apply this to $f(x)=x^n$ for \cref{eq:node-spectral}.
\end{proof}

\subsection{Geometric nodes and the Gaussian basis pair}
\begin{theorem}[The Gaussian Newton pair]
\label{thm:gaussian-newton}
Put $P_k(x)=\prod_{j=0}^{k-1}(x-q^j)$. Then
\begin{align}
 x^n&=\sum_{k=0}^n\G nk P_k(x),
 \label{eq:gaussian-newton}\\
 P_n(x)&=\sum_{k=0}^n(-1)^{n-k}q^{\binom{n-k}2}\G nk x^k.
 \label{eq:gaussian-newton-inverse}
\end{align}
\end{theorem}
\begin{proof}
By \cref{eq:node-h}, the first coefficient is
$h_{n-k}(1,q,\ldots,q^k)$, which is $G_{n,k}(q)$ by
\cref{eq:finite-recip}. The second is the elementary symmetric
coefficient in the same geometric list, and \cref{eq:finite} gives
exactly the stated phase. Equivalently, the first recurrence is
\cref{eq:pascal1} with $a_k=q^k$.
\end{proof}
This is the precise sense in which Gaussian coefficients are
Stirling-type coefficients: the nodes $0,1,2,\ldots$ produce ordinary
Stirling numbers, while $1,q,q^2,\ldots$ produce Gaussian polynomials.
At $q=1$ these geometric nodes all equal one, so
\cref{eq:gaussian-newton} becomes the ordinary binomial expansion
$x^n=\sum_k\binom nk(x-1)^k$, not the ordinary Stirling expansion.

For generic nonzero $q$, the explicit geometric-grid formula is
\begin{equation}
 f[1,q,\ldots,q^k]
 =\frac{q^{-\binom k2}}{(q;q)_k}
   \sum_{j=0}^k(-1)^jq^{\binom{k-j}2}\G kj f(q^j).
 \label{eq:geometric-dd}
\end{equation}
Indeed,
\begin{equation}
 \prod_{i\ne j}(q^j-q^i)
 =(-1)^jq^{jk-j(j+1)/2}(q;q)_j(q;q)_{k-j},
 \label{eq:geometric-denominator}
\end{equation}
obtained by separating $i<j$ from $i>j$. Substitution in
\cref{eq:divided-difference} proves \cref{eq:geometric-dd}.
This is a finite interpolation formula; an infinite Newton expansion
would need separate convergence assumptions.

\subsection{An affine-geometric master connection}
\begin{theorem}[Affine-geometric nodes]
\label{thm:affine}
For nodes $a_j=A+Bq^j$, the connection coefficients are
\begin{equation}
 S_{A+Bq^\bullet}(n,k)
 =\sum_{d=k}^n\binom nd A^{n-d}B^{d-k}\G dk.
 \label{eq:affine-geometric}
\end{equation}
This is a polynomial identity in $A,B,q$, including $B=0$.
\end{theorem}
\begin{proof}
First take $B$ invertible and write $x=A+By$. Then
$P_k^{\boldsymbol a}(x)=B^kP_k(y)$. Expand $(A+By)^n$ ordinarily,
then expand each $y^d$ by \cref{eq:gaussian-newton}. The coefficient of
$B^kP_k(y)$ is \cref{eq:affine-geometric}. Both sides are polynomials,
so the identity descends from $\Q(A,B,q)$ to $\Z[A,B,q]$ and can be
specialized at $B=0$.
\end{proof}
This theorem is a compact engine for several families that otherwise
look unrelated. The next section makes its $q$-arithmetic specializations
explicit.

\section{\texorpdfstring{$q$}{q}-Stirling coefficients and operator normal ordering}
\label{sec:stirling}
\subsection{The convention and three closed formulae}
Define
\begin{equation}
 F_k(x)=\prod_{j=0}^{k-1}(x-\qnum j),\qquad
 x^n=\sum_{k=0}^nS_q(n,k)F_k(x),\qquad
 F_n(x)=\sum_{k=0}^ns_q(n,k)x^k.
 \label{eq:qstirling-def}
\end{equation}
Thus
\begin{equation}
 S_q(n+1,k)=S_q(n,k-1)+\qnum k S_q(n,k).
 \label{eq:qstirling-rec}
\end{equation}
This is the convention used for the type-$A$ array in \cite{dingzeng}.
The rescaled array $q^{\binom k2}S_q(n,k)$ is also frequently called a
$q$-Stirling array; it must not be silently substituted for $S_q$.

\begin{theorem}[Explicit Gaussian formulae for $q$-Stirling arrays]
\label{thm:qstirling}
For $0\le k\le n$,
\begin{align}
 S_q(n,k)&=h_{n-k}(\qnum0,\qnum1,\ldots,\qnum k),
 \label{eq:qstirling-h}\\
 S_q(n,k)&=(1-q)^{k-n}
       \sum_{d=k}^n(-1)^{d-k}\binom nd\G dk,
 \label{eq:qstirling-gaussian}\\
 S_q(n,k)&=\frac{q^{-\binom k2}}{\qfac k}
       \sum_{j=0}^k(-1)^{k-j}q^{\binom{k-j}2}\G kj\qnum j^n.
 \label{eq:qstirling-spectral}
\end{align}
The signed inverse coefficients satisfy
\begin{align}
 s_q(n,k)&=(-1)^{n-k}e_{n-k}(\qnum0,\ldots,\qnum{n-1}),
 \label{eq:qstirling-first}\\
 s_q(n,k)&=(1-q)^{k-n}
 \sum_{d=k}^n(-1)^{d-k}\binom dk
        q^{\binom{n-d}2}\G nd.
 \label{eq:qstirling-first-gaussian}
\end{align}
All apparent singularities in \cref{eq:qstirling-gaussian,eq:qstirling-first-gaussian}
at $q=1$ are removable. Formula \cref{eq:qstirling-spectral} is initially
for generic nonzero $q$; its rational expression equals the same polynomial.
\end{theorem}
\begin{proof}
The symmetric-function formulae are \cref{thm:node}.
Since $\qnum j=(1-q)^{-1}-(1-q)^{-1}q^j$, substitution in
\cref{eq:affine-geometric} proves \cref{eq:qstirling-gaussian}.
For the spectral formula, the pairwise differences are
\[
 \qnum j-\qnum i=
 \begin{cases}q^i\qnum{j-i},&i<j,\\-q^j\qnum{i-j},&i>j.
 \end{cases}
\]
Therefore
\[
 \prod_{i\ne j}(\qnum j-\qnum i)
 =(-1)^{k-j}q^{jk-j(j+1)/2}\qfac j\qfac{k-j}.
\]
Use \cref{eq:node-spectral} and the identity
$jk-j(j+1)/2=\binom k2-\binom{k-j}2$.
For the final formula put $y=1-(1-q)x$, so that
$F_n(x)=(-1)^n(1-q)^{-n}P_n(y)$. Insert
\cref{eq:gaussian-newton-inverse}, expand $y^d$ ordinarily, and compare
$x^k$. Polynomiality from the definitions proves the cancellation claims.
\end{proof}

The first few rows are
\begin{equation}
\begin{array}{c|lllll}
 n\backslash k&0&1&2&3&4\\\hline
 0&1\\
 1&0&1\\
 2&0&1&1\\
 3&0&1&2+q&1\\
 4&0&1&3+3q+q^2&3+2q+q^2&1
\end{array}
\label{eq:qstirling-table}
\end{equation}
In particular $S_q(n,1)=1$ for $n\ge1$. The limit $q\to1$ recovers
ordinary Stirling numbers, while $q=0$ gives
$S_0(n,k)=\binom{n-1}{k-1}$ for $1\le k\le n$: all positive-index
nodes have become one, and \cref{eq:qstirling-h} counts weak
compositions of $n-k$ into $k$ positions.

\subsection{Colored and type-\texorpdfstring{$B$}{B} node specializations}
\begin{corollary}[Nodes $\qnum{r+mj}$]
\label{cor:colored}
Fix integers $m\ge1$ and $r\ge0$, and define $S_q^{(m,r)}(n,k)$ using
the nodes $a_j=[r+mj]_q$. Then
\begin{equation}
 S_q^{(m,r)}(n,k)
 =(1-q)^{k-n}\sum_{d=k}^n
       \binom nd(-q^r)^{d-k}\Gb dk{q^m}.
 \label{eq:colored}
\end{equation}
It is a polynomial with nonnegative integer coefficients. For $(m,r)=(2,1)$
its defining recurrence is
\begin{equation}
 S_q^{(2,1)}(n+1,k)=S_q^{(2,1)}(n,k-1)+[2k+1]_qS_q^{(2,1)}(n,k),
 \label{eq:typeB}
\end{equation}
the type-$B$ convention discussed in \cite{dingzeng}.
\end{corollary}
\begin{proof}
Write $[r+mj]_q=(1-q)^{-1}-q^r(1-q)^{-1}(q^m)^j$ and use
\cref{thm:affine}. Nonnegativity follows from the complete homogeneous
formula \cref{eq:node-h}; \cref{eq:typeB} is \cref{eq:node-Srec}.
\end{proof}
No signed-permutation interpretation is needed to prove
\cref{eq:colored}. It gives a finite Gaussian formula for the entire
node family, including its boundary values $S_q^{(m,r)}(n,0)=[r]_q^n$.

\subsection{The \texorpdfstring{$q$}{q}-Weyl relation and a \texorpdfstring{$q$}{q}-Dobinski identity}
Let $(Xf)(x)=xf(x)$ and
\begin{equation}
 D_qf(x)=\frac{f(x)-f(qx)}{(1-q)x},\qquad D_qx^n=\qnum n x^{n-1}.
 \label{eq:jackson-def}
\end{equation}
On polynomials this is defined by the second equality even at $x=0$.
A direct calculation gives
\begin{equation}
 D_qX-qXD_q=I.
 \label{eq:qweyl}
\end{equation}

\begin{theorem}[Normal ordering]
\label{thm:normal-order}
Put $C_q(n,k)=q^{\binom k2}S_q(n,k)$. Then, as polynomial operators,
\begin{equation}
 (XD_q)^n=\sum_{k=0}^n C_q(n,k)X^kD_q^k.
 \label{eq:normal-order}
\end{equation}
The polynomial $T_{n,q}(z)=\sum_kC_q(n,k)z^k$ has the identity
\begin{equation}
 T_{n,q}(z)=\frac1{e_q(z)}
       \sum_{j\ge0}\qnum j^n\frac{z^j}{\qfac j}.
 \label{eq:qdobinski}
\end{equation}
This holds formally in $z$ over $\Q(q)$; for $0<q<1$ it also holds
analytically when $|(1-q)z|<1$.
\end{theorem}
\begin{proof}
The relation \cref{eq:qweyl} implies
$D_qX^k=\qnum k X^{k-1}+q^kX^kD_q$ by induction on $k$.
Multiplying the proposed normal form by $XD_q$ gives the coefficient
recurrence
$C_q(n+1,k)=\qnum kC_q(n,k)+q^{k-1}C_q(n,k-1)$.
This is exactly \cref{eq:qstirling-rec} after the stated rescaling,
and the initial operators agree.
For \cref{eq:qdobinski}, evaluate \cref{eq:qstirling-def} at $x=\qnum j$.
When $j\ge k$,
\[
 F_k(\qnum j)=q^{\binom k2}\frac{\qfac j}{\qfac{j-k}},
\]
and the value is zero when $j<k$. Multiply by $z^j/\qfac j$, sum over
$j$, and interchange the finite $k$-sum with the formal series. The
result is $e_q(z)T_{n,q}(z)$. Analytically, $\qnum j$ is bounded for
$0<q<1$, so the same sums converge in the stated disk.
\end{proof}
The normalization in \cref{eq:qdobinski} is therefore forced by the
$q$-Weyl operator, not chosen by analogy with the ordinary Dobinski formula.

\section{\texorpdfstring{$q$}{q}-Eulerian polynomials and the Stirling--Eulerian bridge}
\label{sec:eulerian}
For a permutation $\pi=(\pi_1,\ldots,\pi_n)$, let
$\operatorname{Des}(\pi)=\{i:\pi_i>\pi_{i+1}\}$,
$\des(\pi)=|\operatorname{Des}(\pi)|$, and
$\maj(\pi)=\sum_{i\in\operatorname{Des}(\pi)}i$. Define
\begin{equation}
 A_n(t,q)=\sum_{\pi\in\mathfrak S_n}t^{\des(\pi)}q^{\maj(\pi)}
       =\sum_{j=0}^{n-1}A_{n,j}(q)t^j \qquad(n\ge1).
 \label{eq:qeulerian}
\end{equation}
The major index, rather than an unspecified inversion statistic, is
important in the identities below.

\begin{theorem}[$q$-Worpitzky and a finite Eulerian sum]
\label{thm:worpitzky}
For $n\ge1$,
\begin{equation}
 \sum_{m\ge0}[m+1]_q^n t^m=\frac{A_n(t,q)}{(t;q)_{n+1}}.
 \label{eq:qeulerian-gf}
\end{equation}
Consequently, for $m\ge0$,
\begin{align}
 [m+1]_q^n&=\sum_{j=0}^{n-1}A_{n,j}(q)\G{m+n-j}n,
 \label{eq:qworpitzky}\\
 A_{n,j}(q)&=\sum_{r=0}^j(-1)^rq^{\binom r2}\G{n+1}r[j-r+1]_q^n.
 \label{eq:qeulerian-explicit}
\end{align}
\end{theorem}
\begin{proof}
Expand $[m+1]_q^n$ as the weight sum of vectors
$(a_1,\ldots,a_n)\in\{0,\ldots,m\}^n$, with weight $q^{\sum a_i}$.
Sort the labels by decreasing $a_i$, resolving equal entries by
increasing label. This gives a unique permutation $\pi$ such that
$a_{\pi_1}\ge\cdots\ge a_{\pi_n}$, with a strict inequality at every
descent of $\pi$. Put $d=\des(\pi)$ and subtract from the $i$th sorted
entry the number of descents at positions at least $i$. The resulting
sequence is
$ m-d\ge b_1\ge\cdots\ge b_n\ge0$, and its weight is shifted by
$\maj(\pi)$. Conversely, this construction reconstructs exactly the
vectors producing $\pi$.
Write $e_0=m-d-b_1$, $e_i=b_i-b_{i+1}$ for $1\le i<n$, and $e_n=b_n$.
Then $m=d+\sum_{i=0}^ne_i$ and $\sum b_i=\sum_{i=1}^ni e_i$.
The generating function for these gaps is
\[
 t^dq^{\maj(\pi)}\prod_{i=0}^n\frac1{1-tq^i}.
\]
Sum over permutations to get \cref{eq:qeulerian-gf}. Extract $t^m$
using \cref{eq:finite-recip} for \cref{eq:qworpitzky}; multiply by
$(t;q)_{n+1}$ and use \cref{eq:finite} for
\cref{eq:qeulerian-explicit}.
\end{proof}

\begin{theorem}[Gaussian Stirling--Eulerian transforms]
\label{thm:stirling-eulerian}
For $1\le k\le n$,
\begin{equation}
 q^{\binom k2}\qfac k S_q(n,k)
 =\sum_{\ell=1}^k q^{k(k-\ell)}A_{n,\ell-1}(q)
       \G{n-\ell}{k-\ell}.
 \label{eq:stirling-eulerian}
\end{equation}
Conversely,
\begin{equation}
 A_n(t,q)=\sum_{k=1}^n
 q^{\binom k2}\qfac kS_q(n,k)
       t^{k-1}(tq^{k+1};q)_{n-k}.
 \label{eq:eulerian-stirling}
\end{equation}
\end{theorem}
\begin{proof}
Set $U_k=q^{\binom k2}\qfac kS_q(n,k)$. The Newton expansion at
$x=[m+1]_q$ gives
$[m+1]_q^n=\sum_{k=1}^nU_kG_{m+1,k}(q)$.
Since
\[
 \sum_{m\ge0}G_{m+1,k}(q)t^m=\frac{t^{k-1}}{(t;q)_{k+1}},
\]
comparison with \cref{eq:qeulerian-gf} proves
\cref{eq:eulerian-stirling}.
To invert it, we first establish the finite identity
\begin{equation}
 \sum_{r=0}^N\G Nr q^{r(r-1)}u^r(uq^r;q)_{N-r}=1.
 \label{eq:eulerian-kernel}
\end{equation}
The coefficient of $u^j$ on the left equals
\[
 \G Nj q^{\binom j2}
   \sum_{r=0}^j(-1)^{j-r}q^{\binom r2}\G jr,
\]
which is $\delta_{j,0}$ by \cref{cor:annihilator}. Set
$N=n-\ell$ and $u=tq^{\ell+1}$ in \cref{eq:eulerian-kernel} and
multiply by $t^{\ell-1}$. This expresses $t^{\ell-1}$ as
\[
 \sum_{k=\ell}^n q^{k(k-\ell)}\G{n-\ell}{k-\ell}
          t^{k-1}(tq^{k+1};q)_{n-k}.
\]
Insert this expansion for each monomial of $A_n(t,q)$ in
\cref{eq:eulerian-stirling}. These basis polynomials have distinct
lowest powers, so their coefficients are unique. The resulting
coefficient of the $k$th basis polynomial is \cref{eq:stirling-eulerian}.
\end{proof}
Identity \cref{eq:stirling-eulerian} is classical; its normalization and
history are recorded in \cite{dingzeng}. The proof here exhibits it as
another triangular coefficient transform, directly parallel to the
ordinary Stirling--Eulerian relations in the parent monograph.
For example $A_3(t,q)=1+(2q+2q^2)t+q^3t^2$, and at $(n,k)=(3,2)$
both sides of \cref{eq:stirling-eulerian} equal
$2q+3q^2+q^3$.

\section{Bell polynomials and the geometric-alphabet master formula}
\label{sec:bell}
\subsection{The ordinary Bell engine remains ordinary}
We use the complete exponential Bell polynomial
\begin{equation}
 \Y_n(x_1,\ldots,x_n)
 =n!\!\sum_{\substack{m_1,\ldots,m_n\ge0\\\sum j m_j=n}}
       \prod_{j=1}^n\frac1{m_j!}\left(\frac{x_j}{j!}\right)^{m_j},
 \qquad \Y_0=1.
 \label{eq:bell-def}
\end{equation}
Its generating identity is
\begin{equation}
 \exp\!\left(\sum_{j\ge1}x_j\frac{z^j}{j!}\right)
 =\sum_{n\ge0}\Y_n(x_1,\ldots,x_n)\frac{z^n}{n!}.
 \label{eq:bell-egf}
\end{equation}
To prove it, expand each exponential
$\exp(x_jz^j/j!)$ and multiply; the coefficient condition is exactly
$\sum j m_j=n$. The partial Bell polynomial $\Bcal_{n,k}$ is the same
sum with the additional condition $\sum m_j=k$. Thus
$\Y_n=\sum_k\Bcal_{n,k}$. This is the Bell convention of the parent
article, with $\Y_n$ kept distinct from Bernoulli numbers.

\begin{theorem}[A master coefficient formula for finite $q$-products]
\label{thm:bell-product}
Let $N_1,\ldots,N_s\ge0$, and let $a_\nu,\beta_\nu$ belong to the
coefficient algebra. Define powers of unit series by formal logarithm
and exponential. Put
\begin{equation}
 F(z)=\prod_{\nu=1}^s(a_\nu z;q)_{N_\nu}^{-\beta_\nu},\qquad
 L_j=\sum_{\nu=1}^s\beta_\nu a_\nu^j[N_\nu]_{q^j}.
 \label{eq:Lj}
\end{equation}
Then
\begin{align}
 \coef zn F(z)
 &=\frac1{n!}\Y_n(0!L_1,1!L_2,\ldots,(n-1)!L_n),
 \label{eq:bell-product}\\
 &=\sum_{\substack{m_1,\ldots,m_n\ge0\\\sum j m_j=n}}
       \prod_{j=1}^n\frac{L_j^{m_j}}{j^{m_j}m_j!}.
 \label{eq:bell-product-explicit}
\end{align}
Every sum is finite and involves no auxiliary recurrence.
\end{theorem}
\begin{proof}
Expand the logarithm of each factor:
\[
 -\beta_\nu\sum_{r=0}^{N_\nu-1}\log(1-a_\nu q^rz)
 =\sum_{j\ge1}\frac{\beta_\nu a_\nu^j}{j}
      \left(\sum_{r=0}^{N_\nu-1}q^{rj}\right)z^j.
\]
Summing gives $\log F(z)=\sum_{j\ge1}L_jz^j/j$. Apply
\cref{eq:bell-egf}, and then \cref{eq:bell-def}. All rearrangements
are coefficientwise finite because the outer $q$-products are finite.
\end{proof}
The list $1,q,\ldots,q^{N-1}$ is sometimes called a geometric alphabet.
Its power sum is $p_j=[N]_{q^j}$. Logarithms convert multiplication of
factors into addition of these power sums; Bell polynomials convert the
result back to coefficients. No new ``$q$-Bell polynomial'' is necessary
for this operation.

\begin{corollary}[Gaussian coefficients as ordinary Bell polynomials]
\label{cor:gaussian-bell}
With $p_j=[N]_{q^j}$,
\begin{align}
 q^{\binom k2}\G Nk
 &=\frac1{k!}\Y_k(p_1,-1!p_2,2!p_3,\ldots,
                     (-1)^{k-1}(k-1)!p_k),
 \label{eq:gaussian-bell-e}\\
 \G{N+k-1}k
 &=\frac1{k!}\Y_k(p_1,1!p_2,2!p_3,\ldots,(k-1)!p_k),\quad N\ge1.
 \label{eq:gaussian-bell-h}
\end{align}
\end{corollary}
\begin{proof}
Apply \cref{thm:bell-product} to $(-z;q)_N$ and to $1/(z;q)_N$,
then use \cref{thm:finite}.
\end{proof}
For instance,
\begin{equation}
 q\G N2=\frac{p_1^2-p_2}{2},\qquad
 q^3\G N3=\frac{p_1^3-3p_1p_2+2p_3}{6}.
 \label{eq:gaussian-bell-small}
\end{equation}
These identities compress subset sums into power sums. The visible
subtractions are compatible with the nonnegative polynomial model;
they are a change of coordinates, not a loss of positivity.

\subsection{Products, reciprocal products, and mixed parameters}
For a formal parameter $\lambda$, \cref{thm:bell-product} gives
\begin{equation}
 \coef zn(z;q)_N^{-\lambda}
 =\sum_{\sum j m_j=n}\prod_{j=1}^n
       \frac{\lambda^{m_j}[N]_{q^j}^{m_j}}{j^{m_j}m_j!}.
 \label{eq:poch-power}
\end{equation}
It is polynomial in $\lambda$ of degree at most $n$. With $\lambda$ a
positive integer, one may instead expand each factor ordinarily:
\begin{equation}
 \coef zn(z;q)_N^{-\lambda}
 =\sum_{d_0+\cdots+d_{N-1}=n}
       q^{\sum_{j=0}^{N-1}jd_j}
       \prod_{j=0}^{N-1}\binom{\lambda+d_j-1}{d_j}.
 \label{eq:poch-power-factor}
\end{equation}
The proof is multiplication of $(1-q^jz)^{-\lambda}$; hence
\cref{eq:poch-power,eq:poch-power-factor} form an explicit identity
between a partition sum and an ordinary-binomial weighted sum.
More general mixed ratios require only the signed power sums $L_j$ of
\cref{eq:Lj}. In particular, negative $\beta_\nu$ introduce numerator
factors without changing the coefficient mechanism.

\section{Jackson differentiation and geometric Taylor expansion}
\label{sec:jackson}

\subsection{Coefficient extraction and Newton polynomials}
Recall that
\[
 \Dq f(x)=\frac{f(x)-f(qx)}{(1-q)x},\qquad
 \Dq x^n=\qnum n x^{n-1}.
\]
For a formal power series the expression at $x=0$ means its regular
power-series extension. Consequently,
\begin{equation}
 (\Dq^n f)(0)=\qfac n\coef xn f(x).
 \label{eq:q-coef-extraction}
\end{equation}
In this section $q$ is generic; in analytic statements it is enough to
exclude the roots of unity that make the displayed denominators vanish.
The elementary operator formulae below agree with the conventions of
\cite[\S17.2(iv)]{dlmf17}.

Set
\[
 P_k(x;a)=\prod_{j=0}^{k-1}(x-aq^j),\qquad P_0(x;a)=1.
\]
Factoring the $k-1$ common factors in $P_k(x;a)-P_k(qx;a)$ gives
\begin{equation}
 \Dq P_k(x;a)=\qnum k P_{k-1}(x;a).
 \label{eq:q-factorial-lowering}
\end{equation}
Indeed $P_k(qx;a)=q^{k-1}(qx-a)P_{k-1}(x;a)$, whereas
$P_k(x;a)=(x-aq^{k-1})P_{k-1}(x;a)$; their difference is
$(1-q^k)xP_{k-1}(x;a)$.

\begin{theorem}[Polynomial geometric Taylor formula]
For every polynomial $p$ of degree at most $d$,
\begin{equation}
 p(x)=\sum_{k=0}^{d}
       \frac{(\Dq^k p)(a)}{\qfac k}P_k(x;a).
 \label{eq:q-taylor}
\end{equation}
When $a,qa,\ldots,q^ka$ are distinct,
\begin{equation}
 \frac{(\Dq^k p)(a)}{\qfac k}=p[a,qa,\ldots,q^ka].
 \label{eq:q-derivative-dd}
\end{equation}
The same equality holds for arbitrary functions defined on these nodes,
with the Jackson derivative evaluated by its finite-difference formula.
\end{theorem}
\begin{proof}
The monic polynomials $P_k$ are a basis. Apply $\Dq^r$ to an expansion
$p=\sum c_kP_k$ and put $x=a$. Terms with $k<r$ vanish after
differentiation, and terms with $k>r$ contain the factor $P_{k-r}(a;a)=0$.
The remaining term is $c_r\qfac r$. This proves \cref{eq:q-taylor}.
Uniqueness of Newton interpolation proves \cref{eq:q-derivative-dd}.
For arbitrary node data choose their interpolating polynomial: both sides
depend on exactly those data, not on the choice of an extension.
\end{proof}
No convergence assertion for an infinite geometric Taylor series has been
made. Polynomial identities and finite interpolation require no such
assertion.

\subsection{An explicit finite-difference formula}
Let $T_qf(x)=f(qx)$. Since $T_qx^{-1}=q^{-1}x^{-1}T_q$, induction gives
\[
 \Dq^n=\frac{x^{-n}}{(1-q)^n}
             \prod_{r=0}^{n-1}(1-q^{-r}T_q).
\]
Expand this product by the finite $q$-binomial theorem, using reciprocity
of Gaussian coefficients, to obtain
\begin{equation}
 \boxed{\displaystyle
 (\Dq^nf)(x)=\frac{1}{(1-q)^nx^n}
   \sum_{j=0}^n(-1)^j q^{\binom j2-(n-1)j}
             \G nj f(q^jx).}
 \label{eq:q-difference-explicit}
\end{equation}
Here $x\ne0$; for a regular function at zero the limiting or formal
coefficient interpretation is used. The factor $q^{-(n-1)j}$ is
essential. Omitting it is a common consequence of treating multiplication
by $x^{-1}$ as though it commuted with $T_q$.

\subsection{Leibniz and ordered multinomial formulae}
\begin{theorem}[Jackson product rule]
For functions on the required geometric nodes,
\begin{equation}
 \Dq^n(fg)(x)=\sum_{j=0}^n\G nj
       (\Dq^{n-j}f)(q^jx)(\Dq^jg)(x).
 \label{eq:q-leibniz}
\end{equation}
For $r$ factors its iterated form is
\begin{equation}
 \Dq^n\!\left(\prod_{i=1}^rf_i\right)(x)
 =\sum_{n_1+\cdots+n_r=n}
    \frac{\qfac n}{\prod_i\qfac{n_i}}
    \prod_{i=1}^r(\Dq^{n_i}f_i)
              \left(q^{n_{i+1}+\cdots+n_r}x\right).
 \label{eq:q-multi-leibniz}
\end{equation}
\end{theorem}
\begin{proof}
The first-order identity is
$\Dq(fg)=(\Dq f)g+f(qx)\Dq g$.
Also $\Dq(h(q^jx))=q^j(\Dq h)(q^jx)$.
Differentiate the right side of \cref{eq:q-leibniz}. The coefficients of
$(\Dq^{n+1-j}f)(q^jx)(\Dq^jg)(x)$ add to
$q^j\G nj+\G n{j-1}=\G{n+1}j$.
This proves the induction. Repeated application to the first factor and
the remaining product yields \cref{eq:q-multi-leibniz}; the successive
Gaussian coefficients telescope to the displayed multinomial.
\end{proof}

The same coefficient occurs in a genuinely noncommutative binomial
formula. If two elements of an associative algebra satisfy $YX=qXY$,
then
\begin{equation}
 (X+Y)^n=\sum_{k=0}^n\G nk X^kY^{n-k}.
 \label{eq:quantum-binomial}
\end{equation}
Multiplication on the left gives the coefficient
$\G n{k-1}+q^k\G nk$, which is $\G{n+1}k$; this proves the formula.
Its ordering convention is part of its statement. In contrast, for
ordinary commuting $x,y$, the coefficients of $(x+y)^n$ are ordinary
binomial coefficients. Likewise $e_q(x+y)=e_q(x)e_q(y)$ does not hold
for commuting arguments in general.

\section{An exact Jackson chain rule through divided differences}
\label{sec:chain}
The preceding distinctions become especially important for composition.
There is no general identity
$\Dq(f\circ g)=(\Dq f)\circ g\cdot\Dq g$.
For $f(u)=u^2$ and $g(x)=x^2$, its two sides would be
$\qnum4x^3$ and $\qnum2^2x^3$, which differ for generic $q$.
A correct finite chain rule uses divided differences. The general
construction belongs to the divided-difference chain-rule literature
\cite{floaterlyche}; the following proof also explains its relation to
coefficient calculus.

\begin{theorem}[Path form of the divided-difference chain rule]
Let $x_0,\ldots,x_n$ be distinct and put $y_i=g(x_i)$. Suppose first that
the $y_i$ are distinct. For $n\ge1$,
\begin{equation}
 \boxed{\displaystyle
 (f\circ g)[x_0,\ldots,x_n]
 =\sum_{k=1}^n\ \sum_{0=i_0<i_1<\cdots<i_k=n}
   f[y_{i_0},\ldots,y_{i_k}]
   \prod_{r=1}^k g[x_{i_{r-1}},\ldots,x_{i_r}].}
 \label{eq:dd-chain}
\end{equation}
Whenever confluent limits exist, the same formula extends to coinciding
$y_i$ by replacing divided differences by their confluent values.
For $x_i=aq^i$ and $a\ne0$, multiplication of the right side by
$\qfac n$ gives $(\Dq^n(f\circ g))(a)$.
\end{theorem}
\begin{proof}
We give the finite matrix argument in detail. Let $T$ be the
$(n+1)\times(n+1)$ upper bidiagonal matrix whose diagonal is
$x_0,\ldots,x_n$ and whose first superdiagonal consists of ones. Direct
back substitution in $(zI-T)R=I$ shows that
\[
 R_{ij}=\frac1{\prod_{h=i}^j(z-x_h)}\quad(i\le j).
\]
For a polynomial $h$, applying polynomial functional calculus, or
expanding this rational expression by partial fractions, therefore gives
\begin{equation}
 h(T)_{ij}=h[x_i,\ldots,x_j].
 \label{eq:matrix-dd}
\end{equation}
For completeness, the partial-fraction coefficient of $1/(z-x_r)$ is
$\prod_{i\le h\le j,\,h\ne r}(x_r-x_h)^{-1}$; multiplying by $h(x_r)$
and summing is exactly the explicit divided-difference formula.

More generally, let $A$ be upper triangular with diagonal $a_i$ and
strict upper part $N$. The finite expansion
\[
 (zI-A)^{-1}
 =\sum_{k=0}^n\bigl((zI-\operatorname{diag}A)^{-1}N\bigr)^k
                (zI-\operatorname{diag}A)^{-1}
\]
shows that its $(0,n)$ entry is a sum over increasing paths
$0=i_0<\cdots<i_k=n$. The contribution of a path is
\[
 \frac{A_{i_0i_1}\cdots A_{i_{k-1}i_k}}
      {(z-a_{i_0})\cdots(z-a_{i_k})}.
\]
The same partial-fraction calculation gives
\[
 f(A)_{0n}=\sum_{	ext{paths}}
       f[a_{i_0},\ldots,a_{i_k}]
       A_{i_0i_1}\cdots A_{i_{k-1}i_k}.
\]
Take $A=g(T)$, use \cref{eq:matrix-dd}, and use
$f(g(T))=(f\circ g)(T)$. This proves \cref{eq:dd-chain} for
polynomials. For arbitrary data with distinct $y_i$, interpolate $g$
on the $x_i$ and $f$ on the $y_i$; every term is unchanged. Passage to
confluent limits proves the stated extension when the necessary
smoothness or analyticity is available. Finally apply
\cref{eq:q-derivative-dd} to the geometric nodes.
\end{proof}

For $n=2$ the two paths give the particularly transparent identity
\begin{align}
 (f\circ g)[x_0,x_1,x_2]
 &=f[y_0,y_2]g[x_0,x_1,x_2] \notag\\
 &\quad+f[y_0,y_1,y_2]g[x_0,x_1]g[x_1,x_2].
 \label{eq:dd-chain-two}
\end{align}
As $q\to1$, geometric nodes coalesce and divided differences become
ordinary derivatives divided by factorials. Thus this is a finite-node
precursor of ordinary Fa\`a di Bruno, rather than an unjustified
replacement of all factorials by $q$-factorials.

\section{Composition in a \texorpdfstring{$q$}{q}-divided-power normalization}
\label{sec:composition}

\subsection{Products, powers, and logarithms}
Write $Q_n=\qfac n$, with $Q_0=1$, and represent an ordinary commuting
formal series by
\[
 A(z)=\sum_{n\ge0}a_n\frac{z^n}{Q_n}.
\]
This is a change of coefficient normalization over $\Q(q)$, not a new
notion of multiplication or substitution. In particular,
\begin{equation}
 Q_n\coef zn A(z)B(z)=\sum_{k=0}^n\G nk a_kb_{n-k}.
 \label{eq:q-normalized-product}
\end{equation}
For several factors the coefficients are Gaussian multinomials. This
follows immediately by multiplying monomials and using
$Q_n/(Q_kQ_{n-k})=\G nk$.

If $a_0=1$, the ordinary generalized binomial theorem gives the following
closed profile formula. In every profile below $m_j\ge0$,
$\sum_{j=1}^n jm_j=n$, and $M=\sum_jm_j$:
\begin{equation}
 Q_n\coef zn A(z)^\lambda
 =Q_n\sum_{\sum jm_j=n}
      \frac{\fall\lambda M}{\prod_jm_j!}
      \prod_{j=1}^n\left(\frac{a_j}{Q_j}\right)^{m_j}.
 \label{eq:q-normalized-power}
\end{equation}
Here $\fall\lambda M=\lambda(\lambda-1)\cdots(\lambda-M+1)$ is an
\emph{ordinary} falling factorial. For $n\ge1$ the corresponding
logarithm formula is
\begin{equation}
 Q_n\coef zn\log A(z)
 =Q_n\sum_{\sum jm_j=n}
       \frac{(-1)^{M-1}(M-1)!}{\prod_jm_j!}
       \prod_j\left(\frac{a_j}{Q_j}\right)^{m_j}.
 \label{eq:q-normalized-log}
\end{equation}
To prove both identities, put $U=A-1$ and expand
$(1+U)^\lambda=\sum_{M\ge0}\fall\lambda M\,U^M/M!$ or
$\log(1+U)=\sum_{M\ge1}(-1)^{M-1}U^M/M$. The multinomial theorem for
$U^M$ supplies $M!/\prod m_j!$. Since $U$ has positive order, only
finitely many terms contribute to each coefficient.

\subsection{A precise partial Bell normalization}
\begin{definition}
For $G(z)=\sum_{j\ge1}g_jz^j/Q_j$, define
\begin{equation}
 \mathcal C^{(q)}_{n,k}(g_1,g_2,\ldots)
 =\frac{Q_n}{Q_k}\coef zn G(z)^k.
 \label{eq:q-composition-coeff}
\end{equation}
Set $\mathcal C^{(q)}_{0,0}=1$ and
$\mathcal C^{(q)}_{n,0}=0$ for $n>0$.
\end{definition}
This notation fixes one particular $q$-normalization; it does not identify
it with every family called a ``$q$-Bell polynomial'' in the literature.
A direct multinomial expansion yields
\begin{align}
 \mathcal C^{(q)}_{n,k}(g)
 &=\frac{Q_n}{Q_k}
   \sum_{\substack{\sum jm_j=n\\\sum m_j=k}}
       \frac{k!}{\prod_jm_j!}
       \prod_j\left(\frac{g_j}{Q_j}\right)^{m_j}
 \label{eq:q-partial-profile}\\
 &=\frac{Q_n}{Q_k}\frac{k!}{n!}
   \Bcal_{n,k}\left(\frac{1!g_1}{Q_1},
                   \frac{2!g_2}{Q_2},\ldots\right).
 \label{eq:q-partial-ordinary}
\end{align}
If $F(u)=\sum_{k\ge0}f_ku^k/Q_k$, then
\begin{equation}
 Q_n\coef zn F(G(z))
    =\sum_{k=0}^n f_k\mathcal C^{(q)}_{n,k}(g).
 \label{eq:q-composition-law}
\end{equation}
Equations \eqref{eq:q-partial-profile}--\eqref{eq:q-composition-law}
are proved simply by substituting and collecting coefficients. Together
with \cref{eq:q-coef-extraction} they give an at-zero Jackson derivative
formula for formal composition. They do \emph{not} give a shifted
Jackson chain rule at a nonzero point.

There is an instructive warning already at degree three:
\begin{equation}
 \mathcal C^{(q)}_{3,2}(g)
       =\frac{2[3]_q}{[2]_q}g_1g_2.
 \label{eq:q-partial-warning}
\end{equation}
This coefficient is rational, not polynomial, in $q$. Its $q\to1$
limit is $3g_1g_2$, the ordinary partial Bell polynomial. A combinatorial
family with polynomial $q$-weights requires additional structure; it
does not follow from factorial renormalization alone.

\subsection{Substitution matrices and a normalized Riordan group}
For a series $F$ of positive order, let $C(F)$ be the lower triangular
matrix with entries
\[
 C(F)_{n,k}=\frac{Q_n}{Q_k}\coef zn F(z)^k.
\]
Substitution in two stages proves
\begin{equation}
 C(F\circ G)=C(G)C(F).
 \label{eq:q-substitution-matrix}
\end{equation}
The order is important: the inner series occurs in the left matrix.
All matrix entries are finite sums, so the equation is meaningful for
formal series without analytic convergence assumptions.

More generally, set
\begin{equation}
 R(g,f)_{n,k}=\frac{Q_n}{Q_k}\coef zn g(z)f(z)^k,
 \qquad f(0)=0.
 \label{eq:q-riordan}
\end{equation}
When $g(0)$ and $f'(0)$ are units, these matrices form a group and satisfy
\begin{align}
 R(g,f)R(h,\ell)&=R\bigl(g(h\circ f),\ell\circ f\bigr),
 \label{eq:q-riordan-law}\\
 R(g,f)^{-1}&=R\left(\frac1{g\circ\bar f},\bar f\right),
 \label{eq:q-riordan-inverse}
\end{align}
where $\bar f=f^{-1}$ is the ordinary compositional inverse. To prove
the product law, expand $h(u)\ell(u)^k$ in powers of $u$, substitute
$u=f(z)$, and observe that the intermediate $Q_j$ factors cancel.
The inverse then follows by substitution into the product law.
Equivalently, this is diagonal conjugation of the ordinary
power-series Riordan representation by $\operatorname{diag}(Q_n)$.
Thus it extends the parent's matrix viewpoint without introducing a
new associative law.

As a useful check,
\[
 R(e_q(z),z)_{n,k}=\G nk,
\]
and its inverse is $R(E_q(-z),z)$. This recovers Gaussian inversion,
including the factor $(-1)^{n-k}q^{\binom{n-k}2}$.
If polynomials $P_n(x)$ are defined by
$g(z)e_q(xz)=\sum P_n(x)z^n/Q_n$, then
$\Dq P_n(x)=\qnum n P_{n-1}(x)$: apply $\Dq$ in $x$ to the generating
function and compare coefficients. This supplies the corresponding
$q$-Appell lowering property with no ambiguity about normalization.

\section{Ordinary Lagrange inversion with Gaussian input}
\label{sec:lagrange}

\subsection{A closed coefficient theorem}
Let $K$ be a characteristic-zero field containing the needed parameters,
and let $\Phi(w)\in1+wK[[w]]$. There is a unique $g(z)\in zK[[z]]$
satisfying $g=z\Phi(g)$: its coefficient of degree $n$ is fixed by the
previous coefficients. The following formula eliminates that recurrence.

\begin{theorem}[Lagrange--Bell coefficient formula]
For $1\le r\le n$,
\begin{equation}
 \coef zn g(z)^r=\frac rn\coef w{n-r}\Phi(w)^n.
 \label{eq:lagrange}
\end{equation}
In particular, if
\begin{equation}
 \Phi(w)=\prod_{\nu=1}^s(a_\nu w;q)_{N_\nu}^{-\beta_\nu},
 \qquad L_j=\sum_{\nu=1}^s\beta_\nu a_\nu^j[N_\nu]_{q^j},
 \label{eq:lagrange-input}
\end{equation}
then
\begin{align}
 \coef zn g(z)^r
 &=\frac{r}{n(n-r)!}
     \Y_{n-r}\bigl(nL_1,n1!L_2,n2!L_3,\ldots\bigr)
 \label{eq:lagrange-bell}\\
 &=\frac rn
    \sum_{\sum jm_j=n-r}\prod_{j=1}^{n-r}
            \frac{(nL_j/j)^{m_j}}{m_j!}.
 \label{eq:lagrange-profile}
\end{align}
An empty profile contributes one when $n=r$.
\end{theorem}
\begin{proof}
Put $f(w)=w/\Phi(w)$; it is the compositional inverse of $g$. Formal
residue substitution gives
\begin{align*}
 \coef zn g^r
 &=\Res_w w^rf(w)^{-n-1}f'(w)\dd w\\
 &=-\frac1n\Res_w w^r\frac{\dd}{\dd w}f(w)^{-n}\dd w
 =\frac rn\Res_w w^{r-1}f(w)^{-n}\dd w.
\end{align*}
The last equality uses that the residue of a formal derivative is zero.
This is \cref{eq:lagrange}. Residue substitution itself follows by
linearity from monomials: for exponents different from $-1$ the
substituted expression is a derivative, and $f'/f$ has residue one.
Finally, \cref{eq:lagrange-input} gives
$\log\Phi=\sum_{j\ge1}L_jw^j/j$. Apply the proved Bell expansion of an
exponential to $\Phi^n$. This proves both closed forms.
\end{proof}
The prefactor is $r/n$, not $[r]_q/[n]_q$. Ordinary composition and its
formal residue Jacobian have not changed. Replacing this prefactor by a
$q$-ratio generally produces an incorrect inverse.

\begin{example}[A cubic inverse]
The inverse at zero of $f(w)=w(1-w)(1-qw)$ begins
\begin{align*}
 g(z)={}&z+(1+q)z^2+(2+3q+2q^2)z^3\\
 &+(5+10q+10q^2+5q^3)z^4\\
 &+(14+35q+45q^2+35q^3+14q^4)z^5+O(z^6).
\end{align*}
Every coefficient is obtained from \cref{eq:lagrange-profile} using
$L_j=1+q^j$. Equivalently,
$\coef zn g=n^{-1}\sum_{a+b=n-1}\binom{n+a-1}{a}
\binom{n+b-1}{b}q^b$. This is a finite binomial sum, and substitution
into the cubic supplies a separate check.
\end{example}

\subsection{Optional-slot trees and a Gaussian divisibility identity}
A second choice of input reveals polynomial positivity that is not
obvious from Lagrange's factor $1/n$.

\begin{theorem}[Gaussian slot-tree coefficients]
Fix an integer $N\ge1$ and define $T_N(z)$ by
\begin{equation}
 T_N(z)=z\prod_{j=0}^{N-1}(1+q^jT_N(z)).
 \label{eq:slot-tree-eq}
\end{equation}
For $n\ge1$, its coefficient $C_{N,n}(q)=\coef zn T_N(z)$ satisfies
\begin{equation}
 C_{N,n}(q)=\frac1n
   \sum_{\substack{k_1+\cdots+k_n=n-1\\0\le k_i\le N}}
        q^{\sum_{i=1}^n\binom{k_i}2}
        \prod_{i=1}^n\G N{k_i}.
 \label{eq:slot-gaussian}
\end{equation}
Moreover,
\begin{align}
 C_{N,n}(q)&\in\Z_{\ge0}[q],
 \label{eq:slot-positive}\\
 q^{(N-1)(n-1)}C_{N,n}(q^{-1})&=C_{N,n}(q),
 \label{eq:slot-reciprocal}\\
 C_{N,n}(1)&=\frac1n\binom{Nn}{n-1}
 =\frac1{(N-1)n+1}\binom{Nn}{n}.
 \label{eq:slot-fuss}
\end{align}
Thus the numerator of \cref{eq:slot-gaussian} is divisible by $n$
coefficientwise in $\Z[q]$.
\end{theorem}
\begin{proof}
Use Lagrange inversion with $\Phi(w)=(-w;q)_N$, expand each of the $n$
factors by the finite $q$-binomial theorem, and extract degree $n-1$.
This proves \cref{eq:slot-gaussian}.
For integrality and positivity, consider rooted trees in which each
vertex has $N$ labeled child slots $0,\ldots,N-1$, each either empty or
occupied by one subtree. Give an edge in slot $j$ weight $q^j$, and give
each vertex weight $z$. Decomposing at the root gives
\cref{eq:slot-tree-eq}. There are finitely many such trees of each size;
their weight sum is a polynomial with nonnegative integer coefficients.
Uniqueness of the formal solution identifies it with $T_N$.

Reverse every slot label $j\mapsto N-1-j$. A tree of $n$ vertices has
$n-1$ edges, so its exponent $w$ changes to $(N-1)(n-1)-w$. This
weight-reversing bijection proves reciprocity. Finally put $q=1$ in
Lagrange's formula: $\Phi(w)^n=(1+w)^{Nn}$. The equality of the two
ordinary binomial expressions is a direct factorial cancellation.
\end{proof}

For $N=2$ an especially simple form is
\begin{equation}
 C_{2,n}(q)=\frac1n\sum_{j=0}^{n-1}
                \binom nj\binom n{j+1}q^j.
 \label{eq:narayana-slot}
\end{equation}
For example $C_{2,3}=1+3q+q^2$. The formula follows by extracting
$w^{n-1}$ from $(1+w)^n(1+qw)^n$.
These are Narayana polynomials under the displayed normalization.
They are not the same $q$-Catalan family as a Gaussian quotient such as
$\G{2n}n/[n+1]_q$. The name ``$q$-Catalan'' by itself does not specify
a weight statistic or a unique deformation.

\section{Parameter derivatives and weighted Gaussian sums}
\label{sec:derivatives}

\subsection{Differentiating a finite Pochhammer product}
Let $F(a)=(a;q)_N$. Away from its zeros, direct logarithmic
differentiation gives, for $r\ge1$,
\begin{equation}
 L_r(a):=\frac{\dd^r}{\dd a^r}\log F(a)
 =-(r-1)!\sum_{j=0}^{N-1}\frac{q^{jr}}{(1-aq^j)^r}.
 \label{eq:poch-log-derivative}
\end{equation}
Since $F(a+h)/F(a)=\exp(\sum_{r\ge1}L_r(a)h^r/r!)$, the Bell
coefficient identity immediately gives
\begin{equation}
 \frac{\dd^m}{\dd a^m}(a;q)_N
   =(a;q)_N\Y_m(L_1(a),\ldots,L_m(a)).
 \label{eq:poch-derivative-bell}
\end{equation}
The apparent poles cancel on the left; the displayed quotient form is
used only where its factors are defined. For $(a;q)_N^{-\lambda}$,
replace every $L_r$ by $-\lambda L_r$. In this way all derivative orders
are finite profile sums, without repeatedly differentiating a product.

The Euler operator $a\partial_a$ brings the ordinary Stirling numbers
back into the picture:
\begin{align}
 \sum_{k=0}^N(-1)^kq^{\binom k2}\G Nk k^s a^k
   &=(a\partial_a)^s(a;q)_N \notag\\
   &=\sum_{r=0}^s\Sg sr a^r\partial_a^r(a;q)_N.
 \label{eq:weighted-gaussian-derivatives}
\end{align}
For $s=0$ use $0^0=1$. The first equality follows by applying the
operator to the finite $q$-binomial theorem. The second follows from
$(a\partial_a)^s=\sum_r\Sg sr a^r\partial_a^r$, proved either on each
monomial or by the ordinary Stirling recurrence. Combining
\cref{eq:poch-derivative-bell,eq:weighted-gaussian-derivatives} gives
closed formulae for all polynomially weighted signed Gaussian sums.
The Stirling numbers here are ordinary because $a\partial_a$, rather
than a Jackson operator, is being used.

\subsection{Euler derivatives in the base \texorpdfstring{$q$}{q}}
Write $\theta=q\partial_q$ and define the rational functions
\begin{equation}
 R_r(x)=\left(x\frac{\dd}{\dd x}\right)^{r-1}\frac{x}{1-x}
       =\Li_{1-r}(x),\qquad r\ge1.
 \label{eq:Rr}
\end{equation}
The polylogarithm notation is optional: the following finite expression
makes the rational nature explicit:
\begin{equation}
 R_r(x)=\sum_{\ell=1}^r
       (\ell-1)!\Sg r\ell\frac{x^\ell}{(1-x)^\ell}.
 \label{eq:Rr-stirling}
\end{equation}
To verify it, note that for $|x|<1$ the left side is
$\sum_{m\ge1}m^{r-1}x^m$. Alternatively, differentiating the right side
and using the Stirling recurrence proves the same formula from
$R_1=x/(1-x)$; rational continuation removes the restriction $|x|<1$.
For $r\ge2$ the Eulerian form is
\begin{equation}
 R_r(x)=\frac{xA_{r-1}(x,1)}{(1-x)^r},
 \label{eq:Rr-eulerian}
\end{equation}
by the ordinary specialization of the proved Eulerian generating
function. Thus both Stirling and Eulerian bases organize derivatives
of the same $q$-product.

For $0\le k\le n$ define, away from zeros and poles of the factors,
\begin{equation}
 K_r(n,k;q)=\sum_{i=1}^k
       \left(i^rR_r(q^i)-(n-k+i)^rR_r(q^{n-k+i})\right).
 \label{eq:gaussian-log-q-derivative}
\end{equation}
Logarithmic differentiation of the Gaussian product gives
$K_r=\theta^r\log\G nk$, hence
\begin{equation}
 \theta^m\G nk=\G nk\Y_m(K_1,\ldots,K_m).
 \label{eq:gaussian-q-derivative-bell}
\end{equation}
Indeed $q\mapsto qe^t$ converts $\theta$ into $\partial_t$, reducing
this to the ordinary Bell exponential formula. The signed ordinary
first-kind Stirling numbers $s(m,r)$ convert Euler derivatives back to
ordinary ones:
\begin{equation}
 \partial_q^m\G nk
   =q^{-m}\sum_{r=0}^m s(m,r)\theta^r\G nk.
 \label{eq:euler-to-ordinary}
\end{equation}
This follows from $q^m\partial_q^m=\theta(\theta-1)\cdots(\theta-m+1)$.
At $q=1$ the separate terms in \cref{eq:gaussian-log-q-derivative}
are singular, even though the Gaussian polynomial is not. A different
local representation gives exact finite values without cancellation of
infinities.

\section{All Taylor coefficients at \texorpdfstring{$q=1$}{q = 1}}
\label{sec:nearone}

\subsection{Bernoulli cumulants}
Use Bernoulli numbers $B_j$ with convention
\begin{equation}
 \frac{t}{e^t-1}=\sum_{j\ge0}B_j\frac{t^j}{j!},
 \qquad B_1=-\frac12.
 \label{eq:bernoulli-convention}
\end{equation}
This is the standard convention in \cite[\S24.2]{dlmfbernoulli}. Put
\begin{equation}
 d=k(n-k),\qquad
 \Delta_s(n,k)=\sum_{j=1}^n j^s-\sum_{j=1}^k j^s
                             -\sum_{j=1}^{n-k}j^s.
 \label{eq:delta-power}
\end{equation}
Empty sums are zero. We state the result for $n\ge1$; the case $n=0$
is the constant polynomial one.

\begin{theorem}[The exponential-coordinate expansion]
For $0\le k\le n$ and $q=e^t$,
\begin{equation}
 \log\frac{\Gb nk{e^t}}{\binom nk}
   =\frac d2t+
     \sum_{r\ge1}\frac{B_{2r}\Delta_{2r}(n,k)}{2r(2r)!}t^{2r}.
 \label{eq:q-one-log}
\end{equation}
The logarithm is its analytic germ at $t=0$. The series converges for
$|t|<2\pi/n$ and is also a formal identity. Define
\begin{equation}
 \kappa_1=\frac d2,\qquad
 \kappa_{2r}=\frac{B_{2r}}{2r}\Delta_{2r}(n,k),\qquad
 \kappa_{2r+1}=0\quad(r\ge1).
 \label{eq:gaussian-cumulants}
\end{equation}
Then
\begin{equation}
 \Gb nk{e^t}=
   \binom nk\sum_{m\ge0}\Y_m(\kappa_1,\ldots,\kappa_m)
                          \frac{t^m}{m!}.
 \label{eq:q-one-bell}
\end{equation}
In particular, every ordinary derivative at $q=1$ is given by the
finite sum
\begin{equation}
 \boxed{\displaystyle
 \left.\partial_q^m\G nk\right|_{q=1}
    =\binom nk\sum_{j=0}^m s(m,j)
                       \Y_j(\kappa_1,\ldots,\kappa_j).}
 \label{eq:q-one-derivatives}
\end{equation}
\end{theorem}
\begin{proof}
Differentiating and using \cref{eq:bernoulli-convention}, then
integrating with zero constant term, gives
\[
 \log\frac{e^u-1}{u}
    =\frac u2+\sum_{r\ge1}\frac{B_{2r}}{2r(2r)!}u^{2r}.
\]
The nearest nonzero zeros of $e^u-1$ are $\pm2\pi i$, so this local
logarithm is analytic for $|u|<2\pi$. In the Gaussian product cancel
the factors of $t$ and the integer factorials, and add these logarithms
with $u=jt$, with signs prescribed by numerator and denominators.
The linear term is $\Delta_1t/2=dt/2$, giving
\cref{eq:q-one-log}. Exponentiation gives \cref{eq:q-one-bell}.
Finally \cref{eq:euler-to-ordinary} evaluated at one gives
\cref{eq:q-one-derivatives}. Every step is valid as a formal series,
and the preceding disk also gives the asserted analytic validity.
\end{proof}

All quantities in these formulae are finite. For example, if $B_m(x)$
denotes the Bernoulli polynomial defined by
$te^{xt}/(e^t-1)=\sum B_m(x)t^m/m!$, then
\begin{equation}
 \Delta_s(n,k)=
 \frac{B_{s+1}(n+1)-B_{s+1}(k+1)-B_{s+1}(n-k+1)+B_{s+1}(1)}{s+1}.
 \label{eq:delta-faulhaber}
\end{equation}
To prove this power-sum formula, subtract the generating functions at
$x+1$ and $x$ to get
$B_{s+1}(x+1)-B_{s+1}(x)=(s+1)x^s$, then telescope.

\subsection{Area statistics and exact central moments}
Choose uniformly a $k$-subset of $\{0,1,\ldots,n-1\}$ and let
$X$ be the sum of its elements minus $\binom k2$. The subset formula
shows that $\mathbb E e^{tX}=\Gb nk{e^t}/\binom nk$.
Thus \cref{eq:gaussian-cumulants} really is a cumulant formula for this
finite distribution. In particular,
\begin{equation}
 \mathbb E X=\frac d2,\qquad
 \operatorname{Var}(X)=\frac{d(n+1)}{12},\qquad
 \mathbb E(X-d/2)^4
     =3\left(\frac{d(n+1)}{12}\right)^2-\frac{\Delta_4}{120}.
 \label{eq:area-moments}
\end{equation}
Here $\Delta_2=d(n+1)$, $B_2=1/6$, and $B_4=-1/30$.
Every odd centered moment is zero, since the complement-in-the-rectangle
bijection sends $X$ to $d-X$; equivalently the centered logarithm in
\cref{eq:q-one-log} is even. Higher centered moments are obtained by
setting the first Bell argument to zero.

The first two ordinary derivatives illustrate why Euler derivatives
and ordinary derivatives must not be confused:
\begin{align}
 \left.\partial_q\G nk\right|_1&=\binom nk\frac d2,
 \label{eq:gaussian-first-q-one}\\
 \left.\partial_q^2\G nk\right|_1
   &=\binom nk\left(\frac{d^2}{4}
             +\frac{d(n+1)}{12}-\frac d2\right).
 \label{eq:gaussian-second-q-one}
\end{align}
The subtraction of $d/2$ converts the second raw moment into the second
falling-factorial moment. These exact identities require no limiting
probability theorem.

\section{Roots of unity: polynomial evaluation and local jets}
\label{sec:roots}

\subsection{Cyclotomic factorization}
Let $\Phi_\ell(q)$ be the $\ell$th cyclotomic polynomial, and define
\begin{equation}
 \nu_\ell(n,k)=\left\lfloor\frac n\ell\right\rfloor
 -\left\lfloor\frac k\ell\right\rfloor
 -\left\lfloor\frac{n-k}\ell\right\rfloor.
 \label{eq:cyclotomic-exponent}
\end{equation}
Factoring each $1-q^j$ into cyclotomic factors gives
\begin{equation}
 \G nk=\prod_{\ell=2}^{n}\Phi_\ell(q)^{\nu_\ell(n,k)},
 \qquad \nu_\ell(n,k)\in\{0,1\}.
 \label{eq:gaussian-cyclotomic}
\end{equation}
Indeed the multiplicity in $(q;q)_m$ is $\lfloor m/\ell\rfloor$,
and $\lfloor u+v\rfloor-\lfloor u\rfloor-\lfloor v\rfloor$ is zero
or one. The $\ell=1$ exponent is zero. This proves not only regularity
at roots of unity but also that every complex zero of a nonconstant
Gaussian polynomial is simple. The multiplicity of a \emph{particular
primitive root} is one or zero; several different cyclotomic factors
can of course occur.

\subsection{The \texorpdfstring{$q$}{q}-Lucas evaluation}
\begin{theorem}[Gaussian Lucas formula]
Let $\zeta$ be a primitive $\ell$th root of unity, $\ell\ge2$, and write
$n=a\ell+b$, $k=c\ell+d$, where $0\le b,d<\ell$. For $0\le k\le n$,
\begin{equation}
 \boxed{\displaystyle
 \Gb nk\zeta=\binom ac\Gb bd\zeta.}
 \label{eq:q-lucas}
\end{equation}
As usual the right side is zero if $d>b$. Values are taken in the
Gaussian polynomial, not in an uncanceled quotient $0/0$.
\end{theorem}
\begin{proof}
At a primitive root,
\[
 \prod_{j=0}^{\ell-1}(1+\zeta^jz)=1-(-z)^\ell.
\]
Split the $n$ factors into $a$ full blocks and a block of length $b$:
\[
 \prod_{j=0}^{n-1}(1+\zeta^jz)
  =(1-(-z)^\ell)^a\prod_{j=0}^{b-1}(1+\zeta^jz).
\]
The coefficient of $z^{c\ell+d}$ on the right is zero if $d>b$;
otherwise it is
$\binom ac(-1)^{c(\ell+1)}\zeta^{\binom d2}\Gb bd\zeta$.
The coefficient on the left is
$\zeta^{\binom k2}\Gb nk\zeta$. To cancel the phase, observe
\[
 \zeta^{\binom k2-\binom d2}=(-1)^{c(\ell+1)}.
\]
For odd $\ell$ the exponent difference is a multiple of $\ell$; for
even $\ell$ use $\zeta^{\ell/2}=-1$ and the fact that $c\ell+2d-1$
is odd. This proves \cref{eq:q-lucas} including the sign.
\end{proof}
This is a classical identity; see \cite{formichella} for its place
among Gaussian identities and extensions to negative arguments.
At $q=-1$ it becomes
\begin{equation}
 \Gb nk{-1}=
 \begin{cases}
  0,&n\text{ even and }k\text{ odd},\\
  \displaystyle\binom{\lfloor n/2\rfloor}{\lfloor k/2\rfloor},
       &\text{otherwise}.
 \end{cases}
 \label{eq:q-minus-one}
\end{equation}

\subsection{All local coefficients after removing a possible zero}
Fix $\zeta$ primitive of order $\ell\ge2$, let
$\nu=\nu_\ell(n,k)$, and put
\[
 H(q)=\frac{\G nk}{\Phi_\ell(q)^\nu}.
\]
When $\ell>n$, take $\nu=0$. Then $H(\zeta)\ne0$. The following is a
finite algebraic prescription for the full Taylor expansion at $\zeta$:
\begin{align}
 L_r&=(-1)^{r-1}(r-1)!
      \sum_{\substack{2\le d\le n\\d\ne\ell}}\nu_d(n,k)
      \sum_{\substack{\omega\text{ primitive}\\d\text{th root}}}
                         (\zeta-\omega)^{-r},
 \label{eq:root-local-cumulants}\\
 \Gb nk{\zeta+h}
   &=\Phi_\ell(\zeta+h)^\nu H(\zeta)
        \sum_{m\ge0}\Y_m(L_1,\ldots,L_m)\frac{h^m}{m!}.
 \label{eq:root-local-expansion}
\end{align}
To prove this, factor $H$ into the linear factors $q-\omega$ of
\cref{eq:gaussian-cyclotomic}. Its $r$th logarithmic derivative at
$\zeta$ is exactly \cref{eq:root-local-cumulants}. Exponentiate its
Taylor series using Bell polynomials and multiply back the removed
factor. The series is analytic up to the nearest zero of $H$; it is
also an identity of formal Taylor series over the cyclotomic field.
Each prescribed coefficient is a finite sum, even though the full
Taylor series for $H$ written as an exponential is infinite.

If $\nu=1$, the first nonzero coefficient is particularly simple:
\begin{equation}
 \left.\partial_q\G nk\right|_{q=\zeta}
       =\Phi_\ell'(\zeta)H(\zeta).
 \label{eq:root-simple-derivative}
\end{equation}
Higher coefficients are finite convolutions of the polynomial
$\Phi_\ell(\zeta+h)$ with the Bell coefficients in
\cref{eq:root-local-expansion}. Thus root evaluation, multiplicity,
and all local derivatives fit the same coefficient calculus, provided
that a zero is removed before taking its logarithm.

\section{Three distinct asymptotic regimes}
\label{sec:asymptotic}
A formula that is uniform for fixed $q$ need not be uniform as $q\to1$.
This section separates three useful regimes instead of silently mixing
them. Finite coefficient formulae give an exact expansion in the first,
infinite products give convergent exponential corrections in the
second, and Euler--Maclaurin gives a uniform Poincar\'e expansion in the
third. The last calculation extends the parent's asymptotic
coefficient machinery to a Gaussian input.

\subsection{Fixed lower index and fixed base}
For fixed $k\ge0$ and $0<|q|<1$, the finite $q$-binomial theorem gives
\begin{equation}
 \boxed{\displaystyle
 \G nk=\frac1{(q;q)_k}
      \sum_{j=0}^k(-1)^jq^{j(n-k+1)+\binom j2}\G kj,
       \qquad n\ge k.}
 \label{eq:fixed-k-exact}
\end{equation}
This is simply the product
$(q^{n-k+1};q)_k/(q;q)_k$. It is an \emph{exact finite} expansion in
powers of $q^n$, not an infinite asymptotic ansatz. In particular,
\[
 \G nk\longrightarrow\frac1{(q;q)_k}\quad(n\to\infty).
\]
For $k\ge1$ the first correction is
$-q^{n-k+1}[k]_q/(q;q)_k$, with the remaining finite terms explicitly
given by \cref{eq:fixed-k-exact}.

\subsection{Both sides large and fixed base}
Fix $0<|q|<1$, and let $k,n-k\to\infty$. Put
$a=q^{k+1}$, $b=q^{n-k+1}$, and $c=q^{n+1}=ab/q$.
Since $(q;q)_m=(q;q)_\infty/(q^{m+1};q)_\infty$,
\begin{equation}
 \G nk=\frac1{(q;q)_\infty}
              \frac{(a;q)_\infty(b;q)_\infty}{(c;q)_\infty}.
 \label{eq:bulk-fixed-q-product}
\end{equation}
The infinite $q$-binomial theorem therefore gives the exact convergent
expansion
\begin{equation}
 \G nk=\frac1{(q;q)_\infty}
 \sum_{r,s,t\ge0}
   \frac{(-1)^{r+s}q^{\binom r2+\binom s2}a^rb^sc^t}
        {(q;q)_r(q;q)_s(q;q)_t}.
 \label{eq:bulk-fixed-q-triple}
\end{equation}
The two numerator series converge absolutely for every finite $a,b$;
the reciprocal series converges absolutely for $|c|<1$, which holds
here. Hence their product and rearrangement are justified.
The logarithm of the relative correction is
\begin{equation}
 \log\bigl((q;q)_\infty\G nk\bigr)
   =-\sum_{j\ge1}\frac{a^j+b^j-c^j}{j(1-q^j)}.
 \label{eq:bulk-fixed-q-log}
\end{equation}
It follows by expanding the logarithm of each product; the absolute
convergence for $|a|,|b|,|c|<1$ justifies interchanging sums. In
particular,
\begin{equation}
 \G nk=\frac1{(q;q)_\infty}
       \left(1-\frac{a+b}{1-q}
                       +O(|a|^2+|b|^2)\right).
 \label{eq:bulk-fixed-q-leading}
\end{equation}
The constant may depend on the fixed base $q$; the $c$ term is of order
$ab$. Higher corrections can be generated either by the triple sum or
by the ordinary Bell expansion of \cref{eq:bulk-fixed-q-log}. None of
these fixed-base formulae asserts uniformity near $q=1$.

\subsection{The double scaling \texorpdfstring{$q=e^{-\tau/n}$}{q = exp(-tau/n)}}
Let $k=k_n$ be an integer, set $\alpha=k/n$ and $b=1-\alpha$, and
assume
\begin{equation}
 \varepsilon\le\alpha\le1-\varepsilon,\qquad
 0<\tau_0\le\tau\le\tau_1<\infty,
 \qquad q=e^{-\tau/n}.
 \label{eq:double-scaling-domain}
\end{equation}
Thus $\alpha$ denotes the actual ratio $k/n$, including any rounding;
there is no unrecorded floor correction. Define a smooth function on
$[0,1]$ by
\begin{equation}
 h_\tau(x)=\log\frac{1-e^{-\tau x}}{\tau x},\qquad h_\tau(0)=0,
 \label{eq:h-tau}
\end{equation}
and set
\begin{align}
 \mathscr S_\tau(\alpha)
 &=-\alpha\log\alpha-b\log b
      +\int_0^1h_\tau(x)\dd x
      -\int_0^\alpha h_\tau(x)\dd x
      -\int_0^b h_\tau(x)\dd x
 \label{eq:rate-integral}\\
 &=\frac{\Li_2(e^{-\tau})-\Li_2(e^{-\tau\alpha})
                  -\Li_2(e^{-\tau b})+\pi^2/6}{\tau}.
 \label{eq:rate-dilog}
\end{align}
The equality follows from
$\frac{\dd}{\dd x}\Li_2(e^{-\tau x})
=\tau\log(1-e^{-\tau x})$ and elementary integration of
$\log(\tau x)$. One can take \cref{eq:rate-integral} as the definition
without using a special function.

\begin{theorem}[Uniform all-order logarithmic expansion]
Under \cref{eq:double-scaling-domain}, for every fixed $M\ge0$,
\begin{align}
 \log\Gb nk{e^{-\tau/n}}
 ={}&n\mathscr S_\tau(\alpha)-\frac12\log(2\pi n\alpha b)
       +\frac12\bigl(h_\tau(1)-h_\tau(\alpha)-h_\tau(b)\bigr)
       \notag\\
 &+\sum_{r=1}^M\frac{C_{2r-1}(\alpha,\tau)}{n^{2r-1}}
       +O\bigl(n^{-(2M+1)}\bigr),
 \label{eq:double-scaling-log}
\end{align}
with an error uniform on the indicated compact parameter set. The
coefficients are explicit:
\begin{align}
 C_{2r-1}(\alpha,\tau)
 ={}&\frac{B_{2r}}{2r(2r-1)}
            \bigl(1-\alpha^{1-2r}-b^{1-2r}\bigr) \notag\\
 &+\frac{B_{2r}}{(2r)!}
  \bigl(h_\tau^{(2r-1)}(1)-h_\tau^{(2r-1)}(\alpha)
       -h_\tau^{(2r-1)}(b)+h_\tau^{(2r-1)}(0)\bigr).
 \label{eq:double-scaling-C}
\end{align}
There are no even negative powers in this logarithmic expansion.
\end{theorem}
\begin{proof}
The product identity has the exact decomposition
\begin{equation}
 \log\Gb nk{e^{-\tau/n}}
  =\log\binom nk+
      \sum_{j=1}^n h_\tau(j/n)
      -\sum_{j=1}^k h_\tau(j/n)
      -\sum_{j=1}^{n-k}h_\tau(j/n).
 \label{eq:double-scaling-exact-decomposition}
\end{equation}
Indeed $1-e^{-\tau j/n}=(\tau j/n)e^{h_\tau(j/n)}$, and the elementary
factors cancel to the ordinary binomial coefficient.

For an endpoint $u\in\{1,\alpha,b\}$ with $un$ integral,
Euler--Maclaurin on mesh $1/n$ reads
\begin{align*}
 \sum_{j=1}^{un}h_\tau(j/n)
 ={}&n\int_0^u h_\tau(x)\dd x+\frac{h_\tau(u)-h_\tau(0)}2\\
 &+\sum_{r=1}^M\frac{B_{2r}}{(2r)!n^{2r-1}}
       \bigl(h_\tau^{(2r-1)}(u)-h_\tau^{(2r-1)}(0)\bigr)
       +O(n^{-(2M+1)}).
\end{align*}
To justify this error with the displayed truncation, apply the usual
periodic-Bernoulli remainder formula one order farther, then absorb
both the next boundary term and the remainder into $O(n^{-(2M+1)})$.
The derivatives of $h_\tau$ of every fixed order are uniformly bounded
on $[0,1]\times[\tau_0,\tau_1]$, because its apparent singularity at
zero is removable. Hence the bounds are uniform. The underlying
Euler--Maclaurin formula follows by integration by parts successively
with periodic Bernoulli functions; its standard remainder form is
recorded in \cite[\S2.10(i)]{dlmfem}.

Finally the factorial Stirling expansion
\cite[\S5.11]{dlmfgamma} gives
\begin{align*}
 \log\binom nk={}&n(-\alpha\log\alpha-b\log b)
                  -\frac12\log(2\pi n\alpha b)\\
 &+\sum_{r=1}^M\frac{B_{2r}}{2r(2r-1)n^{2r-1}}
       (1-\alpha^{1-2r}-b^{1-2r})+O(n^{-(2M+1)}).
\end{align*}
Its remainder is uniform since $k,n-k\ge\varepsilon n$.
Insert these three Euler--Maclaurin expansions and the Stirling
expansion into \cref{eq:double-scaling-exact-decomposition}; collect
the integral, endpoint, and odd-power terms. This gives precisely
\cref{eq:double-scaling-log,eq:double-scaling-C}.
\end{proof}

\subsection{Explicit derivatives and multiplicative coefficients}
No recursive differentiation is needed in \cref{eq:double-scaling-C}.
For $x>0$ and $s\ge1$,
\begin{equation}
 h_\tau^{(s)}(x)=(-1)^{s-1}
  \left(\tau^s\Li_{1-s}(e^{-\tau x})-\frac{(s-1)!}{x^s}\right).
 \label{eq:h-tau-derivative}
\end{equation}
It follows by differentiating
$\log(1-e^{-\tau x})-\log(\tau x)$ and using
\cref{eq:Rr}. At zero the nonsingular values are
\begin{equation}
 h_\tau'(0)=-\frac\tau2,\quad
 h_\tau^{(2r)}(0)=\frac{B_{2r}\tau^{2r}}{2r},\quad
 h_\tau^{(2r+1)}(0)=0\quad(r\ge1),
 \label{eq:h-tau-zero}
\end{equation}
by the Bernoulli logarithm expansion. Equations
\cref{eq:Rr-stirling,eq:h-tau-derivative,eq:h-tau-zero}
therefore make every $C_j$ a finite expression in rational functions
of $e^{-\tau}$, $e^{-\tau\alpha}$, $e^{-\tau b}$ and the parameters.

Exponentiating the leading terms simplifies the prefactor:
\begin{equation}
 \Gb nk{e^{-\tau/n}}
 \sim e^{n\mathscr S_\tau(\alpha)}
     \sqrt{\frac{\tau(1-e^{-\tau})}
         {2\pi n(1-e^{-\tau\alpha})(1-e^{-\tau b})}}
     \sum_{m\ge0}\frac{H_m(\alpha,\tau)}{n^m},
 \label{eq:double-scaling-multiplicative}
\end{equation}
where $C_{2r}=0$ and
\begin{equation}
 H_m=\frac1{m!}\Y_m(1!C_1,2!C_2,\ldots,m!C_m)
    =\sum_{\sum jm_j=m}\prod_{j=1}^m\frac{C_j^{m_j}}{m_j!}.
 \label{eq:double-scaling-H}
\end{equation}
The first terms are
\[
 1+\frac{C_1}{n}+\frac{C_1^2}{2n^2}
   +\frac{C_3+C_1^3/6}{n^3}
   +\frac{C_1C_3+C_1^4/24}{n^4}+\cdots.
\]
This is an all-order Poincar\'e expansion with the compact-parameter
uniformity stated above, not a claim of convergence of the inverse-$n$
series. Its formal $\tau\to0$ limit recovers ordinary binomial Stirling
asymptotics: $h_\tau\to0$, the rate tends to the entropy, and the
prefactor tends to $(2\pi n\alpha b)^{-1/2}$. Endpoint regimes
$\alpha\to0,1$ require a different uniform analysis and are not included
in the theorem.

\section{Gaussian filters and geometric uniform convolutions}
\label{sec:uniform}
This application returns to the Fabius--Rvachev setting of the source
repository. Its geometric-uniform probability model is developed in
\cite{fabiusrepo}. Here the point is the exact coefficient mechanism:
Bell polynomials make finite-convolution moments polynomial in one
geometric tail parameter, and Gaussian coefficients then remove that
parameter by a finite filter.

\subsection{Moments from finite cumulant formulae}
Let $0<q<1$, let $U_0,U_1,\ldots$ be independent uniform random variables
on $[-1,1]$, and set
\begin{equation}
 X_q=(1-q)\sum_{j\ge0}q^jU_j,\qquad
 X_{q,N}=(1-q)\sum_{j=0}^{N-1}q^jU_j.
 \label{eq:uniform-model}
\end{equation}
The first series converges absolutely, even pointwise for all choices
of the $U_j$, since its absolute value is bounded by one. Independence
and $\mathbb E e^{tU}=\sinh(t)/t$ give
\begin{equation}
 \mathbb E e^{tX_q}
   =\prod_{j\ge0}\frac{\sinh((1-q)q^jt)}{(1-q)q^jt}.
 \label{eq:uniform-mgf}
\end{equation}
The finite products converge uniformly on compact $t$-sets: each
factor is $1+O(q^{2j})$ there, and the tails of $X_q$ also converge
uniformly, justifying the expectation limit directly.

\begin{theorem}[Closed even moments]
Let $\gamma_j$ be the cumulants of $X_q$. Then
\begin{equation}
 \gamma_{2r+1}=0\quad(r\ge0),\qquad
 \gamma_{2r}=\frac{2^{2r}B_{2r}}{2r}
                \frac{(1-q)^{2r}}{1-q^{2r}}
       \quad(r\ge1).
 \label{eq:uniform-cumulants}
\end{equation}
Consequently,
\begin{align}
 \mathbb E X_q^{2d}
   &=\Y_{2d}(0,\gamma_2,0,\gamma_4,\ldots,\gamma_{2d})
 \label{eq:uniform-moments-bell}\\
   &=(2d)!\sum_{\sum_{r=1}^d r m_r=d}
          \prod_{r=1}^d\frac1{m_r!}
                 \left(\frac{\gamma_{2r}}{(2r)!}\right)^{m_r}.
 \label{eq:uniform-moments-profile}
\end{align}
For $X_{q,N}$ replace $\gamma_{2r}$ by
$\gamma_{2r}(1-q^{2rN})$.
\end{theorem}
\begin{proof}
From the Bernoulli generating function,
\[
 \log\frac{\sinh z}{z}
    =\sum_{r\ge1}\frac{2^{2r}B_{2r}}{2r(2r)!}z^{2r}.
\]
For example, it follows by writing
$\sinh z/z=e^{-z}(e^{2z}-1)/(2z)$ and applying the previously proved
logarithm expansion; the linear terms cancel. Sum this identity over
$z=(1-q)q^jt$ in a sufficiently small disk around zero. The geometric
sum is $(1-q^{2r})^{-1}$ for the infinite model, or
$(1-q^{2rN})/(1-q^{2r})$ for the finite model. This proves the cumulant
formula. Exponentiating the cumulant series gives the moment formula
by the Bell identity. Each coefficient uses finitely many cumulants,
so the final formulae have no remaining limiting or recurrence step.
\end{proof}

For $q=1/2$, this centered geometric-uniform sum has density
$\operatorname{up}$ on $[-1,1]$ under the Rvachev normalization of
\cite{fabiusrepo}. Thus, in particular,
\begin{equation}
 \int_{-1}^1x^2\operatorname{up}(x)\dd x=\frac19,
 \qquad
 \int_{-1}^1x^4\operatorname{up}(x)\dd x=\frac{19}{675}.
 \label{eq:up-moment-examples}
\end{equation}
The fourth cumulant is $-2/225$ and the second is $1/9$, so the second
number is $3(1/9)^2-2/225$. More generally every even moment at a
rational base $0<q<1$ is rational by
\cref{eq:uniform-moments-profile}.

\subsection{Exact geometric extrapolation and its residual}
\begin{theorem}[Gaussian annihilation filter]
Fix $d\ge0$ and $0<s<1$, and define
\begin{equation}
 w_{d,j}(s)=
 \frac{(-1)^{d-j}s^{\binom{d-j+1}2}}{(s;s)_d}
             \Gb dj s,\qquad 0\le j\le d.
 \label{eq:filter-weights}
\end{equation}
For every polynomial $P$ of degree at most $d$ and every $x_0$,
\begin{equation}
 P(0)=\sum_{j=0}^d w_{d,j}(s)P(x_0s^j).
 \label{eq:filter-exact}
\end{equation}
More precisely, for each integer $r\ge0$,
\begin{equation}
 \sum_{j=0}^d w_{d,j}(s)s^{rj}=
 \begin{cases}
  1,&r=0,\\
  0,&1\le r\le d,\\
  (-1)^d s^{\binom{d+1}2}\Gb{r-1}d s,&r>d.
 \end{cases}
 \label{eq:filter-residual-monomial}
\end{equation}
Accordingly, for any formal series $F(z)=\sum_{r\ge0}c_rz^r$,
\begin{equation}
 \sum_{j=0}^d w_{d,j}(s)F(s^jz)
  =c_0+(-1)^ds^{\binom{d+1}2}
             \sum_{r>d}\Gb{r-1}d s\,c_rz^r.
 \label{eq:filter-residual-series}
\end{equation}
\end{theorem}
\begin{proof}
The finite $s$-binomial theorem gives the weight generating polynomial
\begin{equation}
 \sum_{j=0}^d w_{d,j}(s)x^j
       =\frac{\prod_{h=1}^d(x-s^h)}{(s;s)_d}.
 \label{eq:filter-generating}
\end{equation}
Indeed the coefficient of $x^j$ in the numerator is
$(-1)^{d-j}s^{(d-j)(d-j+1)/2}\Gb d{d-j}s$.
Evaluate \cref{eq:filter-generating} at $x=1$ and at $x=s^r$.
For $1\le r\le d$ one factor is zero. For $r>d$, factor
$-s^h$ from every numerator factor to obtain
\[
 \frac{\prod_{h=1}^d(s^r-s^h)}{(s;s)_d}
   =(-1)^ds^{d(d+1)/2}\frac{(s^{r-d};s)_d}{(s;s)_d}
   =(-1)^ds^{d(d+1)/2}\Gb{r-1}d s.
\]
This proves \cref{eq:filter-residual-monomial}. Apply it to the
monomials of $P$ or, coefficientwise, of $F$ to prove the remaining
claims. The $d=0$ case is included by the empty-product convention.
\end{proof}

One can also derive the weights by Lagrange interpolation of $P$ at
$x_0,x_0s,\ldots,x_0s^d$, when $x_0\ne0$. The polynomial generating
function proof has the advantage that it simultaneously gives the
exact residual on every higher monomial. For analytic $F$, the last
identity holds wherever its power series converges at the finitely
many evaluation points; no summability convention is needed.

\begin{corollary}[Exact finite recovery of geometric-convolution moments]
For $d,N\ge0$, put $s=q^2$. Then
\begin{equation}
 \boxed{\displaystyle
 \mathbb E X_q^{2d}
       =\sum_{j=0}^d w_{d,j}(q^2)
                   \mathbb E X_{q,N+j}^{2d}.}
 \label{eq:moment-extrapolation}
\end{equation}
Here $X_{q,0}=0$ and the zeroth moment is one, including at zero.
\end{corollary}
\begin{proof}
In \cref{eq:uniform-moments-profile}, the finite-$N$ replacement is
$\gamma_{2r}\mapsto\gamma_{2r}(1-y^r)$ with $y=q^{2N}$.
A profile with $\sum rm_r=d$ therefore has degree at most $d$ in $y$.
Thus $\mathbb E X_{q,N}^{2d}=P_d(s^N)$ for a polynomial $P_d$ of
degree at most $d$, whose constant term is $\mathbb E X_q^{2d}$.
Apply \cref{eq:filter-exact} with $x_0=s^N$.
\end{proof}
For $d=1$ the filter is
$P(0)=(P(sx_0)-sP(x_0))/(1-s)$. The higher filters remove every
geometric error $s^{rN}$ with $1\le r\le d$ at once. This exactness
has been proved for the moment sequence through its polynomial tail
parameter; it does not assert exactness for arbitrary pointwise
samples of a density or for a general asymptotic expansion. That
qualification is essential in applications to spline approximations
and Fabius values.

\section{Verification, implementation, and further extensions}
\label{sec:verification}

\subsection{What was checked}
The accompanying \texttt{verify\_identities.py} is a reproducible test
suite, not a proof assistant. Its core uses only Python's standard
library: integers, rational numbers, finite permutations, and elementary
polynomial arithmetic. It performs \textbf{10,267 exact checks}, all of
which passed in the run recorded in \texttt{verification\_results.txt}.
The checks include small and boundary indices, bases $q=1/2,2/3,2$,
polynomial congruences modulo cyclotomic polynomials, and the
specializations $q=1,-1$. The per-family counts are recorded in the
log rather than hidden behind a single success message.

The comparisons deliberately use different constructions. Gaussian
products are compared with polynomial Pascal recurrences and subset
enumeration; $q$-Stirling finite sums with node recurrences;
$q$-Eulerian coefficients with direct permutation enumeration;
Jackson chain rules with barycentric divided differences; inverse
coefficients with independent series substitution; and moments from
cumulants with direct products of elementary uniform moment-generating
series. Gaussian Lucas evaluations are tested as polynomial
congruences, not by floating-point substitutions near roots of unity.
The local derivative tests at $q=-1$ first remove the zero factor.

Finite testing cannot establish a theorem for arbitrary indices. The
proofs above carry that burden. Conversely, a test failure would expose
at least a transcription, convention, or implementation error even if a
nearby theorem were valid. These complementary roles are particularly
valuable for the powers of $q$ in inversion and for the distinction
between $S_q(n,k)$ and $q^{\binom k2}S_q(n,k)$.

\subsection{A numerical check of the double-scaling order}
An optional final experiment uses \texttt{mpmath} with 70 decimal
digits. For $\tau=1.3$ and $\alpha=1/3$, let $L_0(n)$ be the leading
line of \cref{eq:double-scaling-log}, and define
\[
 R_1(n)=\log\Gb n{n/3}{e^{-1.3/n}}-L_0(n)-C_1/n,
 \qquad R_2(n)=R_1(n)-C_3/n^3.
\]
The actual output, rounded here, is
\begin{center}
\begin{tabular}{@{}rrr@{}}
\toprule
$n$ & $n^3R_1(n)$ & $n^5R_2(n)$\\
\midrule
60  & $0.0815396566910922$ & $-0.1977273873208401$\\
120 & $0.0815808310074465$ & $-0.1979993937823302$\\
240 & $0.0815911422914091$ & $-0.1980676188797450$\\
\midrule
Predicted limit & $C_3=0.0815945809653480$ & $C_5=-0.1980903805565352$\\
\bottomrule
\end{tabular}
\end{center}
The approach to the two predicted constants checks the signs and
orders in the formulae. These are high-precision numerical observations,
not interval-certified inequalities. The uniform error estimate is
proved independently in \cref{sec:asymptotic}.

\subsection{How to use the closed formulae}
For a single coefficient, partition-profile formulae avoid generating
unneeded lower coefficients. For a whole triangular array, a recurrence
or truncated-series multiplication is often faster, even when a closed
formula is available. ``Closed'' and ``computationally cheapest'' are
different properties. In particular, alternating Gaussian sums can
lose precision if evaluated numerically near $q=1$; use polynomial
arithmetic, the Bernoulli expansion of \cref{sec:nearone}, or sufficient
exact precision instead. At roots of unity, use the polynomial,
cyclotomic, or Lucas form before evaluating a rational quotient.

The spectral Newton formula is useful when the nodes are distinct and
few, while the homogeneous-symmetric formula remains regular at
coalescing nodes. The Bell formula is most effective when logarithmic
coefficients are simple. The Gaussian filter is exact only after the
finite-degree dependence on the tail parameter has been established.
These selection rules follow from the proofs, rather than from a
preference for one notation.

\subsection{Directions not claimed as results here}
The framework suggests several further investigations. First, one can
replace one geometric alphabet by several alphabets with different bases.
The Bell-product theorem continues immediately after replacing
$[N]_{q^j}$ by the corresponding sum of geometric power sums, but a
single Gaussian connection array need no longer exist. A useful problem
is to classify which node changes preserve particularly simple inverse
matrices.

Second, ordinary multivariate Lagrange inversion can accept the same
finite $q$-product inputs. The determinant and Jacobian factors remain
ordinary; only their coefficients acquire Gaussian or geometric-power
structure. Working out a systematic multivariate analogue of
\cref{eq:lagrange-profile} would extend the parent article's coefficient
calculus without conflating it with noncommutative substitution.

Third, the regimes $q\to\zeta$ with $n\to\infty$, or
$q=e^{-\tau/n}$ with $k/n\to0$, lie outside the uniform theorem proved
here. The local cyclotomic decomposition and the separate fixed-$k$
formula provide starting points for matched expansions, but no uniform
transition theorem for those regimes is asserted. In particular, the
present inverse-$n$ expansion is not advertised as a complete
exponentially improved transseries.

\section{Conclusion}
The main organizing observation is that Gaussian coefficients are
Newton connection coefficients for geometric nodes. It explains why
$q$-binomial inversion, $q$-Stirling sums, colored-node families, and
geometric divided differences fit the same triangular-array calculus.
Weighted permutation sorting then joins that basis calculus to
$q$-Eulerian polynomials. On the other side, logarithms of finite
$q$-products have geometric power-sum coefficients, so the parent's
ordinary Bell engine applies without alteration. It produces powers,
derivatives, local jets, inverse coefficients, cumulants, and asymptotic
corrections in one uniform language.

The places where the language must change have been made explicit.
Jackson differentiation needs geometric shifts; its chain rule is a
path sum of divided differences. Ordinary substitution keeps the
ordinary Lagrange factor $r/n$. A $q$-factorial normalization alone
need not give polynomial $q$-weights. Root-of-unity evaluation must be
performed after cancellation, and fixed-base asymptotics must be
separated from a moving-base limit. With these distinctions in place,
$q$-binomial coefficients become a particularly effective extension of
combinatorial coefficient calculus rather than a source of ambiguous
formal analogies.

\clearpage
\appendix
\section{A formula-selection guide}
\label{app:guide}
The following guide lists the finite mechanisms most useful for exact
calculations. Each entry points to a formula with all indices and
normalizations specified in the main text.
\begin{longtable}{@{}>{\raggedright\arraybackslash}p{0.32\textwidth}>{\raggedright\arraybackslash}p{0.60\textwidth}@{}}
\toprule
Task & Formula and applicable setting\\
\midrule
\endfirsthead
\toprule
Task & Formula and applicable setting\\
\midrule
\endhead
Evaluate a Gaussian polynomial & Subset or rectangle sums
\eqref{eq:subset}--\eqref{eq:rectangle}; no division, hence valid at every
complex $q$.\\[3pt]
Convolve or invert a Gaussian transform &
Vandermonde \eqref{eq:vandermonde}, multiblock
\eqref{eq:multivandermonde}, inversion \eqref{eq:inversion}.\\[3pt]
Change from geometric to affine-geometric nodes &
The finite mixed ordinary/Gaussian sum \eqref{eq:affine-geometric}.\\[3pt]
Compute $q$-Stirling entries without recurrence &
\eqref{eq:qstirling-gaussian}, \eqref{eq:qstirling-spectral},
\eqref{eq:qstirling-first-gaussian}; polynomial cancellation at exceptional
bases.\\[3pt]
Convert Stirling and Eulerian data &
\eqref{eq:stirling-eulerian} and \eqref{eq:eulerian-stirling}, using the
major-index convention \eqref{eq:qeulerian}.\\[3pt]
Extract a coefficient of mixed $q$-products &
The Bell or finite profile formulae
\eqref{eq:bell-product}--\eqref{eq:bell-product-explicit}.\\[3pt]
Differentiate on a geometric grid &
Finite Jackson sum \eqref{eq:q-difference-explicit}; product
\eqref{eq:q-leibniz}; path chain rule \eqref{eq:dd-chain}.\\[3pt]
Compose $q$-normalized ordinary series &
\eqref{eq:q-partial-profile} and \eqref{eq:q-composition-law}; the partial
coefficients can be rational in $q$.\\[3pt]
Invert an ordinary series with finite $q$-product input &
\eqref{eq:lagrange-bell} or \eqref{eq:lagrange-profile}; the prefactor is
ordinary $r/n$.\\[3pt]
Compute derivatives at $q=1$ &
Bernoulli cumulants \eqref{eq:gaussian-cumulants} and the finite
Stirling--Bell sum \eqref{eq:q-one-derivatives}.\\[3pt]
Evaluate or expand at a root of unity &
Lucas \eqref{eq:q-lucas}; regularized local jet
\eqref{eq:root-local-expansion}.\\[3pt]
Extract moving-base asymptotic coefficients &
\eqref{eq:double-scaling-C} and \eqref{eq:double-scaling-H}, under the
compactness assumptions \eqref{eq:double-scaling-domain}.\\[3pt]
Remove a finite geometric tail polynomial &
Gaussian weights \eqref{eq:filter-weights}, exact residual
\eqref{eq:filter-residual-series}, moment application
\eqref{eq:moment-extrapolation}.\\
\bottomrule
\end{longtable}

\clearpage
\begin{thebibliography}{99}
\raggedright
\addcontentsline{toc}{section}{References}

\bibitem{parent}
V.~Reshetnikov, repository maintainer,
\emph{Combinatorial Coefficient Calculus},
LaTeX monograph in the ProveIt repository, live \texttt{main} source
consulted 4 September 2026.
Directory:\\
\url{https://github.com/VladimirReshetnikov/ProveIt/tree/main/Analysis/FabiusFunction/docs/semi-formalized-research-frontiers/drafts/combinatorial-coefficient-calculus/Combinatorial_Coefficient_Calculus}.
The source file in that directory is
\nolinkurl{Combinatorial_Coefficient_Calculus.tex}.

\bibitem{fabiusrepo}
V.~Reshetnikov, repository maintainer,
\emph{The Fabius Function and Rvachev's Up-Function: Exact arithmetic,
probability, Fourier analysis, and corrected sharp asymptotics},
ProveIt repository, live \texttt{main} source consulted 4 September 2026.
\url{https://github.com/VladimirReshetnikov/ProveIt/tree/main/Analysis/FabiusFunction/docs/Fabius_Function_and_Rvachev_Up}.

\bibitem{voigt}
B.~Voigt,
\emph{A common generalization of binomial coefficients, Stirling numbers
and Gaussian coefficients},
Publications de l'Institut de Recherche Math\'ematique Avanc\'ee,
Strasbourg, 229/S-08 (1984), Actes du 8e S\'eminaire Lotharingien de
Combinatoire, pp.~87--89.
\url{https://www.mat.univie.ac.at/~slc/opapers/s08voigt.pdf}.

\bibitem{dingzeng}
M.-J.~Ding and J.~Zeng,
\emph{Some identities involving $q$-Stirling numbers of the second kind
in type B},
The Electronic Journal of Combinatorics \textbf{31}(1) (2024), Paper P1.36.
\url{https://doi.org/10.37236/12147}.

\bibitem{floaterlyche}
M.~S.~Floater and T.~Lyche,
\emph{Two chain rules for divided differences and Fa\`a di Bruno's formula},
Mathematics of Computation \textbf{76} (2007), 867--877.
\url{https://doi.org/10.1090/S0025-5718-06-01916-8}.

\bibitem{formichella}
S.~Formichella and A.~Straub,
\emph{Gaussian binomial coefficients with negative arguments},
arXiv:1802.02684, version 1, 8 February 2018.
\url{https://arxiv.org/html/1802.02684v1}.

\bibitem{dlmf17}
NIST Digital Library of Mathematical Functions,
\emph{\S17.2: $q$-Calculus}.
Finite and infinite $q$-binomial theorems and Jackson calculus.
\url{https://dlmf.nist.gov/17.2}.

\bibitem{dlmf175}
NIST Digital Library of Mathematical Functions,
\emph{\S17.5: ${}_0\phi_0$, ${}_1\phi_0$, ${}_1\phi_1$ Functions}.
$q$-exponential and $q$-binomial series conventions.
\url{https://dlmf.nist.gov/17.5}.

\bibitem{dlmf176}
NIST Digital Library of Mathematical Functions,
\emph{\S17.6: ${}_2\phi_1$ Function}.
The terminating $q$-Chu--Vandermonde sums.
\url{https://dlmf.nist.gov/17.6}.

\bibitem{dlmfbernoulli}
NIST Digital Library of Mathematical Functions,
\emph{\S24.2: Definitions and Generating Functions}.
Bernoulli numbers and polynomials.
\url{https://dlmf.nist.gov/24.2}.

\bibitem{dlmfem}
NIST Digital Library of Mathematical Functions,
\emph{\S2.10: Sums and Sequences}.
Euler--Maclaurin formula and remainder.
\url{https://dlmf.nist.gov/2.10}.

\bibitem{dlmfgamma}
NIST Digital Library of Mathematical Functions,
\emph{\S5.11: Asymptotic Expansions}.
Logarithmic Stirling expansion and error control.
\url{https://dlmf.nist.gov/5.11}.

\end{thebibliography}
\end{document}
